% !TEX root = ../main.tex
\chapter{논의}
\label{ch:discussion}

\ref{ch:results}장에서는 아비트럼 원(Arbitrum One)에서 잰 측정값을 보고했어요. 이 장에서는 \ref{ch:introduction}장의 각 연구 질문(RQ)과 각 목표(O)를, 그 장에서 밝힌 기준에 맞추어 따져 봐요. 그리고 모든 해석은 그 해석이 바탕을 둔 \ref{ch:results}장의 표를 가리켜요. 따로 밝히지 않으면, 같은 기간(점수를 계산한 바로 그 기간)에 대한 말은 스마트 계약(블록체인 위에서 저절로 실행되는 프로그램)의 매칭 코호트(모든 점수가 똑같이 점수를 매기는 계약 묶음)에 대한 것이에요. 그리고 홀드아웃(나중 기간을 떼어 두고 맞혀 보는 시험)에 대한 말은 스펜더(spender, 허락을 받아 대신 쓰는 쪽) 코호트들에 대한 것이에요.

\dissertationSubheading[sec:findings-rq]{연구 질문별 발견}

\dissertationSubsubheading[sub:finding-rq1]{RQ1: 한 층만 걷는 두 걷기}

RQ1은 허락(allowance, 토큰을 얼마까지 대신 써도 된다는 허락) 화살표(엣지)만 걷는 엔도스랭크(EndorseRank)와, 송금(transfer) 화살표만 걷는 AWP(적응형 가중 페이지랭크)가 같은 계약들을 서로 다른 방식으로 줄 세우는지 물어요. 그리고 각 점수가 같은 때의 허락 점검, 송금 점검과 얼마나 가깝게 맞는지도 물어요. 두 점수는 계약들을 다르게 줄 세워요. 두 점수의 일치를 나타내는 구간(믿을 만한 범위)에는 0도 1도 들어 있지 않아요(표~\ref{tab:tie-sensitivity}). 이 갈라짐은 동점(점수가 똑같은 것)이나, 빠진 계약을 다루는 규칙에서 나오지 않아요. 엔도스랭크는 동점인 계약 쌍을 거의 남기지 않아요. 그리고 그래프(점과 화살표로 이은 그림)에 없는 계약을 외톨이 점(화살표가 하나도 없는 점)으로 남겨 두어도, 일치는 바뀌지 않아요(\ref{sec:rank-divergence}절).

이 갈라짐은 두 그래프를 만드는 방식에서 나와요. 보증(믿고 맡김) 그래프는 마지막 허락(가장 최근 허락)을 기록해요. 그리고 보증 그래프의 화살표는 송금 그래프의 절반도 안 돼요. 송금 그래프는 모든 송금을 기록해요 \parencite{do2023awp}. 두 방법은 다시 출발하는 방식(재시작)도 달라요. 페이지랭크(PageRank)는 정상 분포(오래 걸었을 때 각 점에 머무는 비율)의 몫을, 자기가 받은 화살표에 따라 나누어 놓아요 \parencite{page1999pagerank,langville2006pagerank}. 그래서 활발한 보낸 쪽(송신자)에게서 송금을 많이 받지만 허락은 적게 받은 계약은, AWP에서는 높게, 엔도스랭크에서는 그저 그렇게 줄 세워져요. 송금이 모이는 중심(허브)이 아닌데 많은 주인(허락을 해 주는 쪽)이 허락한 계약은 그 반대예요. 가설(H) H1은 뒷받침돼요. 다만 대부분은 만든 방식 때문에 저절로 그래요. 서로 다른 그래프 위를 서로 다르게 다시 출발하며 걷는 두 걷기가 서로 맞으리라고는 처음부터 예상하지 않았기 때문이에요. 실제로는 한 점수가 다른 점수를 대신할 수 없어요. 빠르다는 까닭으로 엔도스랭크를 고른 운영자는, AWP를 더 빠르게 어림한 값이 아니라 다른 순위를 얻게 돼요. 이 갈라짐은 C-PR(결합 페이지랭크, 두 그래프를 함께 걷는 점수)의 전제(바탕이 되는 생각)이기도 해요. 두 걷기가 계약들을 비슷하게 줄 세운다면, 두 그래프를 함께 걷는 걷기 하나가 어느 쪽에 보탤 것도 거의 없을 거예요.

같은 기간에서는 각 점수가 자기가 걷는 화살표를 따라가요. 송금 위를 걷는 걷기들은 송금 가족(송금에서 나온 비슷한 지표 묶음)과 가장 잘 맞고, 허락 가족과는 훨씬 덜 맞아요. 허락 위를 걷는 걷기들은 허락 가족과 가장 잘 맞아요. 그리고 점수가 자기 가족에서 얻은 구간에는 모두 0이 들어 있지 않아요(표~\ref{tab:alignment-hybrid}).

순환성(자기에게서 나온 것으로 자기를 확인하는 문제)이 H2의 범위를 좁혀요. 한 가족의 모든 대리 지표(대신 재는 값)는 짝이 되는 그래프와 같은 이벤트(계약이 남기는 기록)에서 계산해요. 그래서 이 일치는 \textcite{cronbach1955construct}에서 말하는 뜻의 의도한 구성개념(점수가 재려고 하는 생각) 점검이에요. 외부 시험(바깥 기준으로 하는 시험)은 아니에요. 그리고 그 일부는 만든 방식 때문에 저절로 생겨요. AWP의 화살표는 각 송금에 달러(USD) 금액으로 무게(가중치)를 매겨요. 그래서 AWP는 들어온 송금의 금액 합을, 그 송금의 수를 따라가는 것과 거의 같은 만큼 가깝게 따라가요(표~\ref{tab:alignment-transfer}). 동점도 H2의 범위를 좁히지 않아요. 외톨이 점 규칙은 어떤 계수도 바꾸지 않기 때문이에요(표~\ref{tab:tie-sensitivity}). 점수가 자기 화살표로 만든 점검과 맞는다는 좁은 뜻에서는, H2가 뒷받침돼요. 다만 대부분은 만든 방식 때문에 저절로 그래요. 이 비교들은 하나도 미리 정하지 않았어요. 그리고 엔도스랭크와 AWP 사이의 짝지은 차이는 부록~\ref{ch:appendix-er-awp}에 있어요. 이 차이들은 아무 주장도 담고 있지 않아요.

\dissertationSubsubheading[sub:finding-rq2]{RQ2: 계산 비용}

RQ2는 같은 표본에서 각 점수를 계산하는 데 비용이 얼마나 드는지 물어요. 그리고 두 화살표 집합을 함께 걸으면 비용이 얼마나 더 드는지도 물어요. 엔도스랭크는 시간, 메모리, 반복(한 바퀴) 횟수에서 가장 싼 점수예요. 그리고 송금 층을 걷는 두 혼합 점수는 AWP보다 오래 걸려요(표~\ref{tab:benchmark-runtime}). 그래서 두 화살표 집합을 함께 걷는 것은 어느 한쪽만 걷는 것보다 비용이 더 들어요. 반복 한 번에 드는 시간은 $|E|$(화살표 수)와 $|V|$(점의 수)에 따라 늘어요 \parencite{langville2006pagerank}. 그리고 AWP는 엔도스랭크보다 반복이 더 많이 필요해요. AWP의 시간 비율은 감쇠 계수(화살표를 따라갈 확률) $d=0.85$에서 화살표 수 비율보다 한참 위에 있고, $d$에 따라 움직여요(\ref{sec:benchmark-results}절). 그래서 비용은 각 그래프의 크기와 함께, 그리고 그 걷기에 필요한 반복 횟수와 함께 커져요. 벤치마크(속도 재기)는 이 두 가지를 갈라 보지 않아요.

규모에 따른 변화를 잰 표들은 이 결과를 넓혀 주기보다는 이 결과에 단서를 달아요(표~\ref{tab:benchmark-scaling}). 각 단계는 코호트의 부분 표본이고, 그 계약들에 닿는 화살표를 남겨요. 그래서 단계들은 코호트 그래프보다 큰 그래프를 한 번도 시험하지 않아요(\ref{sub:runtime-scaling}절). 감쇠 계수 점검과 금액 기울기(금액이 커질 때 무게가 얼마나 빨리 커지는지 정하는 값) 점검은 일치도(점수가 대리 지표와 맞는 정도)의 모양을 그대로 둬요(표~\ref{tab:robustness-damping}와~\ref{tab:robustness-slope}). H3은 매칭 코호트에서 뒷받침돼요. 이번에도 대부분은 만든 방식 때문에 저절로 그래요. 두 점수가 같은 풀이 프로그램(solver)을 돌리기 때문이에요. 부분 표본에서는 엔도스랭크가 속도에서 앞서는 것이 모든 단계에서 그대로예요. 하지만 시간 비율이 화살표 수 비율을 따라가는 단계는 하나도 없어요. 각 걷기에 필요한 반복 횟수가 단계마다 다르기 때문이에요. 그러니까 H1--H3은 도구를 설명하는 점검이에요. 결과가 미리 정해지지 않았던 시험은 RQ3--RQ5, 그리고 \ref{sub:finding-neutral}절이에요.

\dissertationSubsubheading[sub:finding-rq3]{RQ3: 시간 홀드아웃과 원시 차수}

RQ3에는 두 부분이 있어요. 첫 부분은, 2026년 3월 말에 동결한(점수를 고정한) 점수가 그 뒤 석 달 동안 스펜더 코호트에서 새로 생긴 주인--스펜더 허락과 새 송금자(새로 송금을 보내온 주소)를 예측하는지(미리 맞히는지) 물어요. 모든 점수가 두 라벨(나중에 채점에 쓰는 정답 값)을 예측해요. 그리고 얼마나 잘 예측하는지는 화살표 종류가 정해요(표~\ref{tab:holdout-spenders}). 새 허락 쌍(새로 생긴 주인--스펜더 허락)에서는 받은 허락 수(허락 입차수)가 앞서요. 안쪽 혼합 점수들(두 끝 사이의 섞는 비율을 쓴 C-PR), 엔도스랭크, 허락 층이 그 뒤를 따라요. 송금만으로 만든 점수들은 더 낮게 머물러요. 그 까닭 가운데 하나는, 이 점수들이 송금 화살표가 없는 스펜더에게 0점을 주기 때문이에요. 그런 스펜더는 여섯 가운데 하나꼴이고, 새 허락 쌍을 더 자주 얻었어요. 나머지 스펜더들에서는, W0(3월 말에 동결한 기간)에서 AWP가 이 라벨을 엔도스랭크보다 위에, 받은 허락 수보다 아래에 줄 세워요. 이것은 짝지은 대비 없이 본 결과예요. 그리고 AWP는 W0에서는 C-PR과 비슷하고, W1(미리 등록해 둔 기간)에서는 C-PR보다 아래예요(표~\ref{tab:holdout-receiving}과~\ref{tab:cpr-contrasts}. 이렇게 좁힌 범위는 라벨을 본 뒤에 골랐어요). 새 송금자에서는 송금 입차수(송금을 보내온 서로 다른 주소 수)가 AWP보다 앞서요. 그리고 허락만으로 만든 점수들이 가장 낮아요. W1도 같은 점수들 사이에서 이 모양을 되풀이해요(표~\ref{tab:fresh-holdout}와~\ref{tab:fresh-posthoc}). 이 증거는 스펜더에 관한 거예요. 트레이더(거래하는 쪽)는 \ref{sub:finding-neutral}절에서 다뤄요.

둘째 부분은 어떤 페이지랭크든 자기가 매끄럽게 다듬는(평활하는) 원시 차수(화살표를 그냥 센 수)를 이기는지 물어요. 그리고 이기는 페이지랭크는 하나도 없어요(표~\ref{tab:holdout-diff}과~\ref{tab:fresh-posthoc}). AWP는 두 기간 모두에서 나중의 새 송금자를 송금 입차수보다 못하게 줄 세우고, 구간도 0보다 아래에 있어요. W0에서는 동결 전에 송금을 받은 적이 있는 스펜더만 보아도 그래요. 그래서 AWP가 주는 0점으로는 이 차이를 설명할 수 없어요(표~\ref{tab:holdout-receiving}). 분석 계획의 고르게 다시 출발하기(균등 재시작)를 쓰는 엔도스랭크는, 두 기간 모두에서 새 허락 쌍을 받은 허락 수보다 못하게 줄 세워요. AWP의 활동 재시작(활발한 주소에서 더 자주 다시 출발하기)을 쓰는 엔도스랭크 변형도 그래요. 다만 그 차이는 더 작아요(\ref{sub:endorserank-restarts}절). 고르게 다시 출발하며 걷는 C-PR의 송금 층은, 송금 입차수에 더 크게 못 미쳐요. 이것은 짝지은 대비 없이 본 결과예요(\ref{sec:holdout-results}절, 표~\ref{tab:fresh-holdout}). 이 라벨들에서 다시 출발 규칙이 바꾸는 것은, 걷기가 자기 화살표를 그냥 센 수에 얼마나 못 미치는지예요. 못 미치는지 아닌지는 바꾸지 않아요. 이미 알려진 결과들을 보면, 그냥 센 수는 이기기 어려울 것으로 예상돼요. 자라나는 네트워크에서는 새 연결이 이미 연결을 많이 가진 점에 더 잘 붙어요 \parencite{barabasi1999scaling,newman2001clustering}. 그래서 차수는 링크 예측(앞으로 생길 연결을 맞히기)의 표준 기준선(비교 기준)이에요 \parencite{libennowell2007link}. 이더리움 거래 네트워크에서도 그래요 \parencite{lin2022tracking}. 그리고 웹처럼 차수 상관(이어진 점들끼리 화살표 수가 서로 닮은 정도)이 약한 곳에서는, 페이지랭크가 평균적으로 들어오는 화살표 수(입차수)에 비례해요 \parencite{fortunato2008indegree}. 대부분의 스펜더에게 엔도스랭크는 무게를 다시 단 허락한 주인 수(서로 다른 주인 수)예요(\ref{sub:theory-closed-form}절). 이것이 첫 부분의 답을 뒤집지는 않아요. 그냥 센 수들도 같은 화살표에서 낸 통계이기 때문이에요. 결합 걷기도 둘째 부분에 더 나은 답을 내지 못해요. C-PR은 어느 라벨도 같은 화살표 종류의 그냥 센 수보다 위에 줄 세우지 못해요. 두 기간 모두에서, 새 허락 쌍에서는 받은 허락 수보다 아래에 있고, 새 송금자에서는 송금 입차수보다 아래에 있어요(표~\ref{tab:cpr-contrasts}). 이 대비들은 모두 라벨을 본 뒤에 계산했어요.

라벨은 점수를 매긴 기간 밖에 있어요. 하지만 점수와 같은 구성개념에 속해요. 그리고 어느 라벨에도 이익이나 빚 갚기(상환)의 요소는 하나도 들어 있지 않아요. 그래서 RQ3은 각 구성개념의 시간 예측 타당도(나중 일을 얼마나 잘 미리 맞히는지)를 보여 줘요. 이때 원시 차수가, 페이지랭크가 이겨야 하는 기준선이에요. 그리고 RQ3은 신용 위험은 다루지 않아요. 신용 위험에 가장 가까운 것은 W1의 청산(손실이 너무 커져 거래가 강제로 끝나는 일) 라벨이에요(\ref{sub:fresh-holdout}절). \ref{sub:finding-neutral}절에서 이 라벨로 화살표 종류들을 견줘요.

\dissertationSubsubheading[sub:finding-rq4]{RQ4: 결합 걷기를 그 층들과 견주기}

RQ4와 RQ5는 이 논문의 방법에 관한 질문이에요. 계획이 정한 결합 방식인 C-PR에 대해, 두 질문은 모두 부정적인 답을 얻어요. 각 질문은 자기 규칙에 따라 답을 얻어요. 여기서는 \ref{sub:finding-neutral}절의 R1--R3처럼, 이 규칙들을 고친 점수(등록 뒤에 고쳐서 다시 계산한 점수)에 재분석으로 적용해요. 등록 분석(미리 등록한 그대로 한 분석)도 두 규칙에서 같은 판정에 이르렀어요(부록~\ref{sec:map-registered}). RQ4는 두 층을 함께 걷는 걷기 하나가, 어느 한 층만 걷는 것보다 계약들을 더 잘 줄 세우는지 물어요. 분석 계획에서 정한 규칙 A에 따르면, 그렇지 않아요. W0에서 C-PR은 새 허락 쌍에서 허락 층보다 못하지 않지만, 새 송금자에서는 송금 층에 뒤져요(표~\ref{tab:holdout-diff}). 그리고 같은 기간에서는 보조 규칙이 세 가족 모두에서 실패해요(\ref{sec:coupled-results}절). 이 규칙은 결합 점수가 하나라도 생기기 전에 정했어요. 그래서 섞는 비율 $\lambda$를 이 규칙에 맞추어 조정했을 수는 없어요. 다만 고친 C-PR의 달러 무게, 허락 상한(이보다 큰 금액은 이 값으로 자름), 점수 기간(점수를 계산하는 데 쓰는 기간)은 등록 분석의 결과를 안 뒤에 정했어요. 이 결과는 가장 먼저 떠오르는 다음 단계에 한계를 그어요. 계획에서 정한 방식은 이동 행렬(전이 행렬) 수준에서 두 화살표 종류를 볼록 혼합(더하면 1이 되는 비율로 섞기)하는 거예요. 이 방식은 새 허락 쌍에서는 허락 층보다 앞섰지만, 새 송금자에서는 송금 층에 뒤졌어요.

빼 보기 실험(한 부분을 빼거나 바꾸어 보는 실험)과 층의 무게가 이 답을 설명해요. S-PR(보증 씨앗 페이지랭크)은 송금 걷기를 그대로 두지만, 허락 걷기를 바탕으로 다시 출발하게 해요. S-PR은 흐름 라벨을 송금 걷기만큼, 또는 조금 더 잘 줄 세워요. 하지만 보증 신호는 대부분 잃어요. 걷기가 어디로 갈지를 송금 화살표가 정하면, 보증에서 끌어낸 순간 이동(다시 출발할 곳으로 훌쩍 건너뛰기) 벡터는 별로 바꾸는 것이 없기 때문이에요. 그리고 S-PR은 송금 층의 점들에만 점수를 매겨요. 그래서 그 손실의 일부는 점수를 매기는 범위(커버리지) 때문이에요. C-PR에서는 허락 화살표가 걷기의 일부를 실어 날라요. 그래서 스펜더 홀드아웃에서 허락 신호가 살아남아요. 모든 안쪽 혼합 점수는 새 허락 쌍에서 허락 층보다 위에 있고, 새 송금자에서는 송금 층보다 아래에 있어요(표~\ref{tab:holdout-spenders}와~\ref{tab:holdout-diff}). C-PR의 층들은 허락에 둔 상한 말고는 달러 금액을 한계 없이 그대로 실어요. 그리고 고르게 다시 출발해요. 반면에 AWP는 각 송금의 무게에 한계를 두고, 활발한 보낸 쪽에서 다시 출발해요. AWP는 두 기간 모두에서 새 송금자를 C-PR보다 더 잘 줄 세워요. 그리고 C-PR의 송금 층은 혼자 걸어도 이미 새 송금자를 AWP보다 못하게 줄 세워요(표~\ref{tab:holdout-spenders}와~\ref{tab:fresh-holdout}). 그래서 그 차이의 일부는 결합이 아니라 무게와 다시 출발 방식에서 나와요. 다른 결합 방법은 이 논문에서 평가하지 않아요(\ref{sec:limitations}절).

\dissertationSubsubheading[sub:finding-rq5]{RQ5: 송금 걷기에 더한 허락 화살표}

RQ5는 이 걷기가 송금만 걷는 걷기보다 나은지 물어요. W0에서 한 탐색적(미리 정하지 않고 살펴본) 비교는 새 허락 쌍에서 그렇다고 말했어요. 하지만 그 비교는 라벨을 본 뒤에 한 거예요. 이 질문의 답은 규칙 B가 정해요. 규칙 B는 이 연구가 7월부터 9월까지의 로그(기록)를 뽑기 전에 등록했어요. 그리고 규칙 B의 답은 “아니요”예요. C-PR은 새 허락 쌍에서 다시 송금 층을 이겼어요. 하지만 비열등성(정한 폭보다 더 나쁘지 않음)은 새 송금자에서도, 트레이더의 청산 없이 닫은 비율(거래를 청산 없이 끝낸 비율)에서도 보이지 못했어요. 그리고 이 규칙에는 세 조건이 모두 필요했어요(\ref{sub:fresh-results}절). 첫 번째 민감도 분석(조건을 바꿔 결과가 얼마나 움직이는지 보기)에서는 AWP를 비교 대상으로 삼아요. 거기서는 F1이 성립하고, F2와 F3은 실패해요. C-PR은 새 허락 쌍을 더 높게, 새 송금자를 더 낮게 줄 세워요. 그리고 청산 없이 닫은 비율에서는 둘이 갈리지 않아요(표~\ref{tab:fresh-contrasts}). 이 규칙은 이 연구가 라벨을 뽑기 전에 등록했고, 세 조건 가운데 두 조건에서 실패해요. W1이 확인해 주는 것은 더 좁아요. C-PR의 송금 걷기에 허락 화살표를 더하면, 새 자료에서 새 허락 쌍을 더 잘 줄 세워요. 그 일부는 송금 걷기가 0점을 주는 스펜더들 덕분이에요. 그리고 그 대가로, 새 송금자에서는 등록 문서의 허용 폭(이만큼까지는 나빠도 괜찮다고 정한 폭)보다 더 크게 잃어요. 나머지 스펜더들에서 C-PR이 자기 송금 층을 이기는지는 계산하지 않았어요.

규칙과 따로, 표~\ref{tab:cpr-contrasts}의 대비들은 그 정보가 어디서 오고 어떤 대가를 치르는지 보여 줘요. C-PR은 두 기간 모두에서 두 스펜더 라벨을 엔도스랭크보다 위에 줄 세워요. 자기 허락 층보다 위에 줄 세우는 것과 마찬가지예요. 그래서 송금도 함께 걸으면, 모든 스펜더 라벨에서 허락만 걷는 걷기보다 나아져요. 가장 많이 나아지는 것은 나중에 들어오는 흐름이에요. AWP와 견주면, 나중의 허락에서 얻는 이득은 대부분 송금 화살표가 없어서 송금 걷기가 0점을 주는 스펜더들에게서 나와요. 송금 걷기가 줄 세울 수 있는 스펜더들 사이에서는, 나중의 허락에서 W0에서는 C-PR과 AWP가 비슷하고 W1에서는 C-PR이 앞서요. 반면에 나중의 송금자는 두 기간 모두에서 AWP가 더 잘 줄 세워요. 그러니까 기간 밖에서 허락 층이 송금 걷기에 보태는 것은 주로 점수를 매기는 범위예요. 곧 송금 걷기가 동점으로 남겨 두는 계약들 사이의 순서예요. 그리고 W1에서는 나머지 계약들 사이에서도 나중의 허락에서 작은 이득이 있어요. 활동 재시작을 쓴 엔도스랭크는 허락해 주는 주인에서 다시 출발하는 허락 걷기예요. 이 변형은 나중의 허락에서 W0에서는 C-PR과 비슷하고, W1에서는 C-PR보다 더 잘 줄 세워요. C-PR은 나중의 송금자에서 두 기간 모두 이 변형보다 앞서요. 그리고 기간 안에서는 송금 가족과 시빌 안정성(가짜 주소를 많이 만들어 속이는 시빌 공격을 감안한 안정성) 가족에서 앞서요(\ref{sub:cpr-contrasts}절). 이 연구의 어떤 페이지랭크도, 결합 페이지랭크까지 포함해서, 나중에 새로 들어온 것을 같은 화살표 종류의 그냥 센 수보다 위에 줄 세우지 못해요. 이 대비들은 라벨을 본 뒤에 계산했고, 아무것도 판정하지 않아요.

\dissertationSubsubheading[sub:finding-neutral]{어느 화살표 종류로도 만들지 않은 결과 값에서 허락과 송금 견주기}

RQ3--RQ5의 스펜더 라벨은 자기와 같은 화살표 종류로 만든 점수에 유리해요. 이와 달리 GMX(무기한 선물 거래 서비스) 트레이더의 결과 값은 어느 화살표 종류로도 만들지 않았고, 점수를 동결한 뒤에 쟀어요. 이 결과 값에서 두 화살표 종류를 견주는 비교(\ref{sub:neutral-label-protocol}절)는 새 연구 질문을 보태지 않아요. 이 비교는 주장~C6으로 보고해요(\ref{sub:neutral-label-results}절).

이 비교는 거래 상대 위험(거래하는 상대 때문에 손해를 볼 위험)과 관련이 있어요. 이 논문을 하게 된 까닭이 바로 이 쓰임이에요(\ref{ch:introduction}장). 계약 계정에 레버리지(빌린 돈을 보태 거래하게 해 주는 것)나 신용을 주는 서비스(프로토콜)에는, 자주 강제로 청산되는 거래 상대가 더 위험한 상대예요. 그리고 이 자료에서 채무 불이행(빚을 못 갚는 일)에 가장 가까운 결과가 강제 청산이에요. 주된 기간인 W1에서, 나중에 청산을 피하는 일에 대해 동결일의 허락한 주인 수와 송금 수 사이의 차이 P를 보면, 그 본페로니 97.5\% 구간(여러 번 비교할 때 넓힌 신뢰구간)이 0보다 위에 있어요(표~\ref{tab:neutral-label-contrasts}). 닫은 횟수로 묶음을 나눈(층화한) 구간도 그래요. 그래서 R1이 충족돼요. 등록 문서에서 정한 말로 쓰면, 나중에 청산을 피하는 데에 대해 그냥 센 수의 수준에서 허락이 송금보다 앞서요. 보조 기간인 W0에서는 P가 양수였지만, 어느 쪽이 나은지 가려지지 않았어요. 송금 수도 이 결과 값을 줄 세워요(표~\ref{tab:fresh-holdout}). 그리고 송금 점수들은 활동이 많을수록 커지는 결과 값, 곧 이익을 낸 닫기 수와 실현한 이익(실제로 거둔 이익)도 줄 세워요(표~\ref{tab:fresh-traders}). 송금 점수들은 계약이 얼마나 많이 거래하는지를 기록해요.

네 가지 한계가 이 해석의 범위를 좁혀요. 첫째, R2는 충족되지 않아요. 같은 조건끼리 견준 페이지랭크 짝, 곧 AWP의 다시 출발 규칙을 쓴 엔도스랭크와 AWP는 어느 기간에서도 가려지지 않아요. 그리고 엔도스랭크는 이번에도 자기가 매끄럽게 다듬는 허락한 주인 수보다 아래로 떨어져요(표~\ref{tab:neutral-label-contrasts}). 둘째, 이 비교는 기간마다 허락을 받은 트레이더 몇십 개에 기대고 있어요. 이 트레이더들은 모두 청산을 피했어요. 나머지 트레이더들은 넷 가운데 셋쯤이 청산을 피했어요(표~\ref{tab:neutral-label-recipients}). 그리고 모든 구간이 넓어요. 셋째, 모든 트레이더가 계약이에요. 그래서 계정 종류로는 이 모양을 설명할 수 없어요. 하지만 계약이 맡은 역할로는 설명될 수도 있어요. 허락을 받은 트레이더는 거의 모두 한 계약에게서 허락을 받았어요(\ref{sub:neutral-label-protocol}절). 그래서 이들은 한 무리를 이뤄요. 그리고 이들은 포지션(거래)을 더 많이 닫아요. 이것은 대신 운용해 주거나 돈을 모아 함께 거래하는 계약들의 한 가족과 잘 맞아요. 넷째, W1에서는 R3이 적용돼요. 같은 신호가 이익을 낸 닫기의 비율에서는 더 낮아요. 그래서 이 신호가 줄 세우는 구성개념은 거래를 잘하는 것이 아니라 강제 청산을 피하는 것이에요. 이 신호에서 손실이 아닌 비율이 더 낮은 것은, 일부는 만든 방식 때문에 저절로 생겨요. 처리 과정(파이프라인)이 청산된 닫기를 손실이 아닌 것으로 세기 때문이에요.

\dissertationSubheading[sec:objectives-achievement]{연구 목표를 이룬 정도}

\ref{ch:introduction}장은 다섯 목표 하나하나에 측정과 기준을 붙였어요. 표~\ref{tab:objectives-achievement}은 그것들을 증거와 맞대어 점검해요.

\begin{table}[htbp]
\centering
\caption{연구 목표, \ref{ch:introduction}장에서 밝힌 기준, 그리고 결과.}
\label{tab:objectives-achievement}
\fitwidth{%
\fontsize{7.5}{9}\selectfont%
\begin{tabular}{p{0.05\linewidth}p{0.31\linewidth}p{0.46\linewidth}p{0.13\linewidth}}
\toprule
목표 & 기준 & 증거 & 결과 \\
\midrule
O1 & EndorseRank와 AWP 사이 켄달의 $\tau$에 대한 95\% 구간에 $0$도 $1$도 들어 있지 않아야 해요. 어떤 점수의 구간이 한 가족에서 $0$을 벗어나면, 그 점수는 그 구성개념과 맞는다고 봐요. 그리고 점수가 자기 화살표의 구성개념과 맞는지 보는 것은 의도한 구성개념 점검이에요 & 표~\ref{tab:tie-sensitivity}의 마지막 행: 구간 $[0.181,0.221]$은 $(0,1)$ 안에 완전히 들어 있어요. 표~\ref{tab:alignment-hybrid}: 송금만 걷는 모든 걷기의 송금 가족 값은 $0.700$ $[0.691,0.709]$부터 $0.736$ $[0.728,0.744]$까지예요. 허락을 걷는 두 걷기의 허락 가족 값은 $0.560$ $[0.548,0.571]$와 $0.561$ $[0.550,0.572]$예요. 모든 구간이 $0$을 벗어나요. 대리 지표는 그래프와 같은 이벤트에서 나와요 & 달성(대부분 만든 방식 때문에 저절로) \\
O2 & 하나로 정해 둔 절차 아래에서, 실행 시간 비율을 화살표 수 비율과 함께 적어요 & 표~\ref{tab:benchmark-runtime}: 전체 그래프에서 평균, 표준편차(SD, 값이 흩어진 정도), 메모리, 반복 횟수, $|V|$, $|E|$를 적었어요. 표~\ref{tab:benchmark-scaling}: 단계마다 평균 시간, $|V|$, $|E|$를 적었어요. 시간 비율은 $3.21$이고, 그 옆의 화살표 수 비율은 $2.21$이에요. 감쇠 계수 값에 따라서는 $2.56$부터 $3.21$까지예요(표~\ref{tab:robustness-damping}). C-PR과 S-PR은 AWP보다 느려요. 모든 점수에 풀이 프로그램 하나를 써요 & 달성(대부분 만든 방식 때문에 저절로) \\
O3 & 어떤 점수의 구간에 $0$이 들어 있지 않으면, 그 점수는 라벨을 예측한다고 봐요. 페이지랭크는 짝지은 차이의 구간이 $0$보다 위에 있을 때만 차수보다 낫다고 봐요 & 표~\ref{tab:holdout-spenders}, \ref{tab:fresh-holdout}와~\ref{tab:fresh-posthoc}: 모든 점수가 W0과 W1에서 두 라벨을 모두 예측해요. 표~\ref{tab:holdout-diff}과~\ref{tab:fresh-posthoc}: AWP 빼기 송금 입차수는 $-0.038$ $[-0.043,-0.032]$(W0), $-0.041$ $[-0.047,-0.034]$(W1)이에요. EndorseRank 빼기 받은 허락 수는 $-0.090$ $[-0.100,-0.079]$(W0), $-0.082$ $[-0.090,-0.073]$(W1)이에요. 활동 재시작을 쓴 EndorseRank 빼기 받은 허락 수는 $-0.060$ $[-0.070,-0.052]$(W0), $-0.046$ $[-0.055,-0.038]$(W1)이에요. 표~\ref{tab:cpr-contrasts}: 새 허락 쌍에서 C-PR 빼기 받은 허락 수는 $-0.063$ $[-0.076,-0.052]$(W0), $-0.065$ $[-0.077,-0.054]$(W1)이에요. 새 송금자에서 C-PR 빼기 송금 입차수는 $-0.111$ $[-0.121,-0.100]$, $-0.115$ $[-0.128,-0.105]$예요. 모두 라벨을 본 뒤에 계산했어요 & 첫째 조항은 달성; 둘째 조항은 어떤 페이지랭크도 달성하지 못함 \\
O4 & 2026년 10월 9일의 규칙 A: W0에서 C-PR은 새 허락 쌍에서 C-PR ($\lambda=1$)보다 못하지 않고, 새 송금자에서 C-PR ($\lambda=0$)보다 나아야 해요. 같은 기간에서는 송금 가족과 허락 가족에서 그 점 추정값(값 하나로 낸 추정)이 앞서는 층의 구간 안에 있어야 하고, 시빌 안정성에서는 두 층을 모두 넘어야 해요 & 고친 점수, 재분석. 표~\ref{tab:holdout-diff}: 못하지 않다는 조건은 맞아요($+0.031$ $[0.021,0.042]$). 더 낫다는 조건은 실패해요($-0.030$ $[-0.040,-0.020]$). 같은 기간: 송금($0.400$)은 C-PR ($\lambda=0$)의 구간 밖에 있고, 허락($0.449$)은 C-PR ($\lambda=1$)의 구간 밖에 있어요(표~\ref{tab:alignment-hybrid}). 시빌 안정성은 두 층 사이에 있어요(표~\ref{tab:tau-diff-hybrid}). 세 부분 모두 실패해요(\ref{sec:coupled-results}절) & 미달성(등록 분석에서도 같음, 표~\ref{tab:registered-vs-revised}) \\
O5 & 2026년 10월 9일에 등록한 규칙 B: W1에서 C-PR은 새 허락 쌍에서 C-PR ($\lambda=0$)을 이겨야 해요. 그리고 새 송금자와 트레이더의 청산 없이 닫은 비율에서 0.01보다 더 많이 나빠서는 안 돼요(95\% 구간의 아래 끝이 $-0.01$보다 위) & 고친 점수, 재분석. 표~\ref{tab:fresh-contrasts}: 새 허락 쌍 $+0.170$ $[0.149,0.189]$는 맞아요. 새 송금자 $-0.025$ $[-0.038,-0.015]$는 구간 전체가 $-0.01$보다 아래라서 실패해요. 청산 없이 닫은 비율 $+0.000$ $[-0.010,0.010]$은 아래 끝이 $-0.01$보다 아래라서 실패해요. 민감도 분석 1(AWP가 비교 대상): $+0.159$ $[0.139,0.178]$는 맞고, $-0.075$ $[-0.085,-0.065]$와 $+0.019$ $[-0.026,0.064]$는 실패해요. 사후(결과를 본 뒤에 한) 비교는 표~\ref{tab:cpr-contrasts}에 있어요 & 미달성(등록 분석에서도 같음, 표~\ref{tab:registered-vs-revised}) \\
\bottomrule
\end{tabular}
}
\end{table}

목표 밖의 발견 여덟 가지도 결론을 만드는 데 영향을 줘요. 다섯 가지는 주장을 제한해요. 수정(등록 뒤에 분석을 고친 일)에서 고른 금액 기울기는 같은 기간에서만 점검했어요(\ref{sub:threat-construct}절). 부분 표본 단계들은 코호트 그래프보다 큰 그래프에서 한 번도 돌리지 않았어요(\ref{sub:finding-rq2}절). 트레이더 코호트에서는 어떤 점수도 이익을 낸 닫기의 비율을, 0보다 위에 있는 구간으로 줄 세우지 못해요(\ref{sub:trader-results}절). 두 스펜더 라벨에서 엔도스랭크의 점 추정값은 C-PR의 허락 층보다 조금 위에 있어요. 짝지은 대비는 계산하지 않았어요. 그래서 허락에 무게를 다는 방식이 중요할 수도 있어요(\ref{sub:threat-construct}절). 그리고 새 허락 쌍에서 C-PR이 AWP보다 앞서는 것은 대부분 AWP가 0점을 주는 스펜더들에게서 나와요(\ref{sub:cpr-contrasts}절). 세 가지는 주장에 보태요. C-PR은 두 기간 모두에서 두 스펜더 라벨을 엔도스랭크보다 위에 줄 세워요(\ref{sub:cpr-contrasts}절). 허락을 받은 적이 있는 계약 트레이더는 두 기간 모두에서 하나도 빠짐없이 청산을 피했어요. 그 허락은 거의 모두 한 계약에게서 왔고, 표본은 작아요(\ref{sub:finding-neutral}절). 그리고 깊이가 1인(주인들이 스스로는 허락을 받지 않는) 모든 스펜더에서, 엔도스랭크는 닫힌 식(반복 계산 없이 바로 쓸 수 있는 식)과 맞아요(\ref{sub:theory-closed-form}절).

\dissertationSubheading[sec:theoretical-implications]{이론에 주는 뜻}

이 결과들은 이론의 세 분야와 관련이 있어요. 금융 그래프에서 퍼뜨림(전파)이 무엇을 뜻하는지, 온체인(블록체인 위의) 점수의 구성 타당도(점수가 재려는 것을 정말 재는지), 그리고 가명성(진짜 이름 대신 주소로만 활동함) 아래에서 믿음의 경제학이에요.

\dissertationSubsubheading[sub:theory-propagation]{송금을 넘어선 퍼뜨림의 뜻}

페이지랭크는 처음에 하이퍼링크(웹 페이지끼리 잇는 링크)를 위해 만들어졌어요. 이때 하이퍼링크는 인용 네트워크(서로 가리키고 인용하는 관계)처럼 읽혔어요 \parencite{page1999pagerank}. \textcite{do2023hodgerank}에서는 이 읽는 방식을 이더리움 거래로 옮겨 왔어요. 그리고 AWP는 금액 무게와 나이(얼마나 오래되었는지) 무게, 그리고 활발한 보낸 쪽에서의 다시 출발로 이것을 더 다듬었어요 \parencite{do2023awp}. 엔도스랭크는 이것을 ERC-20(토큰 공통 규칙) 허락에 적용해요. 허락에서는 스펜더가 자기 대신 움직이게 해 주는 주인이, 쓸 권한에 대해 방향이 있는 보증을 내줘요 \parencite{vogelsteller2015erc20}. RQ1의 갈라짐은, 화살표가 뜻하는 것과 걷기가 다시 출발하는 곳을 바꾸면 같은 점들 위에서도 순위가 바뀐다는 것을 보여 줘요. 그래서 그래프에 바탕을 둔 평판은 방법들의 가족이에요. 각 방법은 자기 화살표가 담는 관계에 비추어 타당해요. P2P(컴퓨터끼리 직접 주고받는) 믿음 시스템과 웹 믿음 시스템은 반대쪽에서 출발해 같은 결론에 이르렀어요. 아이겐트러스트(EigenTrust)에서는 씨앗 집합(믿을 만하다고 미리 정한 점들)과 화살표의 정의가 결과를 거의 정해요 \parencite{kamvar2003eigentrust}. 트러스트랭크(TrustRank)에서도 마찬가지예요 \parencite{gyongyi2004trustrank}. 그 방법들의 화살표는 메시지를 주고받는 겹친 네트워크(오버레이)나 하이퍼링크 그래프에 속했어요. 이와 달리 엔도스랭크의 화살표는 토큰 표준 자체에서 나와요. C-PR은 한 걸음 더 나아가, 걷기 하나가 같은 주소들의 두 관계를 모두 따라가게 해요. 그래서 계약의 점수는 그 계약이 받는 돈과 그 계약이 받은 쓸 권한에 한꺼번에 기댈 수 있어요.

홀드아웃은 퍼뜨림이라는 비유가 가리는 단서 하나를 보태요. 화살표 종류를 바꾸면 점수가 미리 내다볼 수 있는 것이 바뀌었어요. 두 기간 모두에서, 새 허락 쌍에서는 허락 화살표로 만든 점수가 앞섰어요. 그 일부는 송금을 걷는 걷기들이 0점을 주는 스펜더들 덕분이었어요. 그리고 새 송금자에서는 송금을 걷는 걷기들이 앞섰어요. 기간 밖에서는 어떤 걷기도 자기 화살표를 그냥 센 수에 무언가를 보태지 못했어요. AWP와 C-PR의 층들, 그리고 다시 출발 규칙이 어느 쪽이든 엔도스랭크는 모두 자기 라벨을 원시 차수보다 못하게 줄 세웠어요(\ref{sub:finding-rq3}절). 이 자료에서는 어떤 화살표를 고르느냐가, 점수가 어떤 스펜더들을 줄 세울 수 있는지 정해요. 그리고 그 화살표를 그냥 센 수가 그 스펜더들을 가장 잘 줄 세워요. 두 층이 한 걷기를 함께 쓸 때, 허락 화살표가 걷기를 직접 이끄는 곳에서는 허락 신호가 살아남아요. 허락 화살표가 다시 출발의 씨앗만 정하는 곳에서는 그 신호가 대부분 사라져요(\ref{sub:finding-rq4}절). 그리고 결합 걷기는 모든 스펜더 라벨을 허락만 걷는 걷기보다 위에 줄 세워요. 하지만 나중의 송금자에서는 송금만 걷는 걷기보다 아래에 머물러요(\ref{sub:finding-rq5}절).

\dissertationSubsubheading[sub:theory-closed-form]{깊이가 1인 스펜더: 걷기가 그냥 센 수에 보태는 것}

허락 그래프들은 코호트 계약이 해 주거나 받은, 다섯 토큰의 허락만 담아요(\ref{sec:graph-construction}절). 그래서 점들은 대부분 나가는 화살표만 있는 순수한 주인이거나, 들어오는 화살표만 있는 순수한 스펜더예요. 코호트 계약은 둘 다일 수 있어요. 그래서 그래프에는 한 걸음보다 긴 길도 있어요. 주인들이 허락을 하나도 받지 않은 스펜더에서는, 고르게 다시 출발하는 걷기를 닫힌 식으로 쓸 수 있어요.

\emph{명제(참이라고 보이려는 말).} 이동 행렬 $P$를 쓰는 페이지랭크를 식~\eqref{eq:pagerank}로, 감쇠 계수 $d$와 고르게 다시 출발하기로 계산한다고 해요. 그러면 매달린 점(나가는 화살표가 없는 점)에 고인 몫도 고르게 퍼져요. 그때 들어오는 화살표가 없는 점은 모두 같은 점수 $c$를 가져요. 그리고 들어오는 이웃이 모두 들어오는 화살표가 없는 점이면, 그런 점 $v$의 점수는 다음과 같아요.
\begin{equation}
\label{eq:closed-form}
r_v = c\,\bigl(1 + d\,\delta(v)\bigr), \qquad \delta(v)=\sum_{u} P_{uv},
\end{equation}
곧 $v$로 들어오는 이동 확률의 합이에요. 엔도스랭크에서는 $P_{uv}=w(u,v)/o(u)$예요. 그래서 $\delta(v)$는 $v$의 주인 정규화 입강도(주인마다 나누어 맞춘, 들어오는 무게의 합)예요. 다시 말해, 각 주인의 보증은 그 주인이 허락하는 모든 스펜더에게 나뉘어 들어가요. C-PR에서 $P_{uv}$는 식~\eqref{eq:coupled-transition}의 결합 이동 확률이에요. 그리고 여기서 들어오는 화살표가 없는 이웃이란, 두 층 어디에서도 들어오는 화살표가 없는 이웃이에요.

\emph{유도(식을 끌어내는 과정).} 고정점(계산을 한 번 더 해도 바뀌지 않는 값)에서는 $r_v = (1-d+d\,r_D)/|V| + d\sum_u r_u\,P_{uv}$예요. 여기서 $r_D$는 매달린 점들이 가진 몫이에요. 들어오는 화살표가 없는 점은 첫째 항만 받아요. 그래서 그런 점은 모두 점수 $c=(1-d+d\,r_D)/|V|$를 가져요. 이것은 \ref{sub:sybil-bound}절의 상수 $c$예요. $v$의 들어오는 이웃 $u$가 모두 그런 점이면, $r_u=c$를 넣어서 식~\eqref{eq:closed-form}을 얻어요.

$T_{\mathrm{obs}}$의 엔도스랭크 그래프에서, 코호트 스펜더 6{,}224개 가운데 5{,}232개는 깊이가 1이에요. 이 논문에서 계산한 엔도스랭크는 그 하나하나에서 식~\eqref{eq:closed-form}과 1\% 안으로 맞아요. 더 깊은 스펜더는 스스로 허락을 받는 주인을 가지고 있어요. 이 가운데서는 21.9\%가 그렇게 맞아요. 모든 스펜더를 통틀면 그 비율은 87.5\%예요. 그리고 동결일 $t_1$의 W0 스펜더를 통틀면 70.8\%예요. 코호트 스펜더에서 엔도스랭크와 $\delta$ 사이의 켄달의 타우($\tau$, 두 줄 세우기가 얼마나 같은 순서인지 나타내는 수)는 $0.908$이에요. 이것은 점 추정값이에요(스크립트 \texttt{run\_same\_window.py}와 \texttt{run\_holdout.py}). 이 비율들은 C-PR의 합친 그래프(두 층을 합친 그래프)에서는 계산하지 않았어요. 그 그래프에서는 송금도 받은 주인에게 들어오는 화살표가 있어요. 그래서 그 주인의 스펜더들은 깊이가 1이 아니에요.

여기서 두 가지 결과가 따라 나와요. 첫째, 깊이가 1인 스펜더에서 엔도스랭크는 허락한 주인 수에 무게만 다시 달아요. 많은 스펜더를 허락하는 주인은 스펜더 하나하나에서 더 적게 세요. 시간 무게와 금액 무게 말고는, 이 정규화(나누어 맞추기)가 엔도스랭크를 그냥 센 수와 다르게 만드는 것의 전부예요. 그리고 새 허락 쌍은 그냥 센 수가 더 잘 줄 세워요(\ref{sub:finding-rq3}절). 그러니까 이런 스펜더에서는, 엔도스랭크에 대한 RQ3 둘째 부분의 부정적인 답에 분석적인 근거(식으로 따져 본 근거)가 있어요. 걷기가 보탤, 퍼뜨려 온 정보가 없었어요. 활동 재시작을 쓰면, 매달린 점의 몫도 다시 출발을 따라가요. 그래서 깊이가 1인 스펜더에서 이 변형은, 주인들의 다시 출발 몫으로 무게를 단, 주인들에 대한 한 걸음짜리 합이에요. 이 변형이 그냥 센 수와 다른 것은 주인들에게 무게를 다는 방식뿐이고, 더 긴 길 때문이 아니에요. 그리고 이 변형도 새 허락 쌍을 그냥 센 수보다 못하게 줄 세워요. 둘째, 그래프에는 정말로 더 깊은 길이 있어요. 그리고 더 깊은 스펜더 대부분에서 점수는 닫힌 식에서 벗어나요. 그래서 이 그래프에는 여러 걸음에 걸친 퍼뜨림이 있어요. 하지만 모든 스펜더를 통틀어서는, 엔도스랭크가 기간 밖에서 그냥 센 수를 이기지 못했어요. 그리고 더 깊은 스펜더만 따로 떼어 시험하지는 않았어요. 에고 네트워크(코호트 계약에 닿는 화살표만 모은 그래프)는 더 긴 길을 잘라 내요(\ref{sub:threat-external}절). 그래서 퍼뜨림을 더 시험하려면, 코호트에서 두 걸음 안의 이웃이나 체인 전체의 허락 그래프 같은 더 깊은 표본이 필요해요. 또는 미리 확인한 계약에서 다시 출발하게 해야 해요.

\dissertationSubsubheading[sub:theory-construct]{여러 구성개념에 걸친 분야 일치도}

점검들은 송금 중심성(송금 흐름의 한가운데에 있는 정도)과 허락 위임(대신 쓰도록 맡기기)을 갈라 보도록 설계했어요. 그리고 각 점수는 자기가 걷는 화살표의 구성개념과 잘 맞고, 혼합 점수는 자기 두 층 사이에 놓여요. 고전적인 구성 타당도 이론은, 서로 다른 구성개념을 재는 도구라면 이런 모양이 나올 거라고 예측해요 \parencite{cronbach1955construct}. 그리고 \textcite{messick1995validity}의 통합된 관점은 한 걸음 더 나아가, 점수를 쓴 결과도 그 타당도의 일부로 쳐요. 같은 이론은 이런 모양이 무엇을 가려 주지 못하는지도 말해 줘요. 두 도구가 서로 다른 구성개념을 재면, 어떤 점검에서 앞서는 쪽은 그 점검을 위해 만든 도구예요. 그래서 계수로 도구들의 순서를 매겨도, 어느 쪽이 평판을 더 잘 재는지는 알 수 없어요. 어느 화살표 종류로도 만들지 않은 결과 값은 예외예요. 어느 도구도 그것을 위해 만들지 않았기 때문이에요. 그래서 이 논문은 거기서 화살표 종류들을 견줘요(\ref{sub:finding-neutral}절). 그런데 거기서도 어느 화살표 종류가 앞서는지는 결과 값에 따라 달라요.

\dissertationSubsubheading[sub:theory-pseudonymity]{값싼 가명 아래의 평판}

공개 장부(누구나 볼 수 있는 블록체인 기록)에서는 신원을 만드는 데 돈이 거의 들지 않아요 \parencite{nakamoto2008bitcoin}. 그리고 이더리움의 계정 방식에서는 새 주소를 만들 때 등록하는 단계가 아예 없어요 \parencite{wood2025yellowpaper}. 그래서 들어오는 송금이나 허락에 보상을 주는 평판 시스템은 모두 시빌 전략에 드러나 있어요. 시빌 전략에서는 공격자가 자기가 움직이는 주소들 사이에 화살표를 만들어요 \parencite{douceur2002sybil}. \ref{sec:sybil-cost-model}절은 허락 그래프를 공격하는 비용을 식으로 정리해요. 허락 그래프에서는 서로 허락해 주는 주소들의 닫힌 농장(공격자가 만든 가짜 주소 무리)이 토큰 잔고를 하나도 필요로 하지 않아요. 이 점이 송금 자전 거래(자기끼리 사고팔아 거래가 많은 척하기)와 달라요. 그리고 가스(거래 수수료) 말고는 돈이 거의 들지 않아요.

\dissertationSubheading[sec:practical-implications]{실제에 주는 뜻과 쓰는 장면}

서비스와 기반 시설(인프라) 회사는 송금 중심지 찾기, 허락을 살피는 평판, 그리고 둘을 합친 걷기 가운데에서 골라야 해요. 잰 결과는 점수 하나를 권하기보다, 쓰는 장면에 따라 점수를 고르는 쪽을 뒷받침해요.

\dissertationSubsubheading[sub:deploy-hybrid]{결합 걷기와 함께 쓰기}

두 신호를 합칠지는 실제로 쓰는 일에 관한 질문이면서 방법에 관한 질문이에요. 그리고 증거는 이 두 질문에 서로 다르게 답해요. 이 자료에서, 두 화살표 종류를 한 걷기 안에 접어 넣는 것은 발표된 송금 걷기에 견주어 도움이 되지 않았어요. C-PR은 AWP보다 비용이 더 들고, 엔도스랭크와 AWP를 한 번씩 푸는 비용을 합친 것보다도 더 들어요(표~\ref{tab:benchmark-runtime}). 그리고 두 기간 모두에서 새 송금자에서 AWP에 뒤져요. AWP가 줄 세울 수 있는 계약들 사이에서는, 새 허락 쌍에서 W1에서만, 그것도 조금만 앞서요(\ref{sub:finding-rq5}절과 \ref{sub:cpr-contrasts}절). 그래서 마음대로 고를 수 있는 운영자라면 두 점수를 따로 두어야 해요. 허락을 살피는 감시에는 엔도스랭크나 허락한 주인 수를 돌리고, 송금 중심지 감사에는 AWP를 돌리는 거예요. 대신 두 번 푸는 비용이 들어요. 달러 금액으로 무게를 단 송금 걷기 하나에 묶인 운영자도 있어요. 그런 운영자에게는, 허락 화살표를 보태면 새 자료에서도 버티는 허락 정보를 얻었어요. 주로 송금 걷기가 0점에 남겨 두는 계약들을 줄 세운 덕분이었어요. 하지만 그 대가로, 들어오는 흐름에서 $0.01$보다 더 크게 잃었어요. 그래서 규칙 B는 충족되지 않아요.

\dissertationSubsubheading[sub:deploy-transfer]{송금 중심지와 흐름 지켜보기}

지금의 흐름을 그려 내는 서비스가 있어요(마켓 메이커(늘 사고팔 값을 내놓는 거래자)와 유동성 공급자(거래할 토큰을 맡겨 두는 사람)의 송금 중심지 찾기, 경로 분석, 송금에 바탕을 둔 규정 준수 어림 규칙(대강 맞는 규칙) 등). 그런 서비스는 그 흐름을 입차수와 받은 양의 합에서 바로 읽을 수 있어요. AWP는 지금 들어오는 흐름과 잘 맞아요(표~\ref{tab:alignment-hybrid}) \parencite{do2023awp}. 그리고 기간 밖에서는, 두 기간 모두에서 나중의 새 송금자를 다른 모든 페이지랭크보다 위에, 하지만 원시 송금 입차수보다는 아래에 줄 세웠어요(표~\ref{tab:holdout-spenders}와~\ref{tab:fresh-posthoc}). 그래서 이 자료에서는, 흐름을 미리 내다보고 싶은 서비스라도 그냥 센 수 대신 AWP를 쓸 까닭이 없어요.

\dissertationSubsubheading[sub:deploy-allowance]{허락을 살피는 쓰임과 빠르게 새로 계산하기}

분명히 밝힌 쓰기 허락이 가장 중요한 곳이 있어요(토큰 금고, 허락이 있어야만 열리는 연결, 허락을 믿음의 행동으로 보는 위험 계산 엔진 등). 그런 곳에서는 서비스가 허락 화살표를 쓸 까닭이 있어요. 가장 싼 꼴은 스펜더의 받은 허락 수예요. 받은 허락 수에는 풀이 프로그램이 필요 없어요. 그리고 두 기간 모두에서 새 허락을 엔도스랭크보다, 그리고 이 연구의 다른 모든 페이지랭크보다 더 잘 예측했어요(표~\ref{tab:holdout-diff}과~\ref{tab:fresh-posthoc}). 그 수를 퍼뜨린 버전인 엔도스랭크는 허락 점검과 가장 잘 맞아요(표~\ref{tab:alignment-hybrid}). 그리고 AWP보다 빨리 풀려요(표~\ref{tab:benchmark-runtime}). 하지만 대부분의 스펜더에서 엔도스랭크는 보증하는 주인들을, 각 주인이 스펜더를 몇 개나 허락하는지로만 비춰요(\ref{sub:theory-closed-form}절). 이 성질이 어떤 라벨에서든 더 낫게 만든다는 것은 보이지 않았어요. 그리고 고른 순간 이동 아래에서는, 엔도스랭크가 닫힌 농장 때문에 싸게 부풀려질 수 있어요(\ref{sec:sybil-cost-model}절). 활동 재시작을 쓴 변형은, 라벨을 본 뒤의 비교에서 그냥 센 수에 더 가까워졌어요. 하지만 두 기간 모두에서 그 수보다 아래에 머물렀어요. 그래서 이 자료에서는 그냥 센 수보다 나은, 퍼뜨린 허락 점수가 없어요.

\dissertationSubsubheading[sub:deploy-industry]{산업에서 쓰는 장면}

아래 장면들은 각 점수가 누구에게 쓸모 있고 누구에게 쓸모없는지 말해 줘요.

\begin{itemize}
\item DEX(탈중앙 거래소) 운영자. 사용자마다 수수료 등급이나 감시 우선순위를 정하는 서비스에는 사용자에 대한 점수가 필요해요. 그런데 그 사용자들은 대부분 외부 소유 계정(개인 키로 움직이는 계정)이고, 이 논문은 외부 소유 계정에 점수를 매기지 않아요(\ref{sec:data-sources}절). 대신 이 점수들은 그 DEX가 연결해 쓰는 라우터(거래 길을 찾아 주는 계약), 금고, 풀(토큰을 모아 둔 곳)을 줄 세울 수 있어요.
\item 대출 서비스와 증거금 거래 서비스. 대출 시장은 담보(빌리는 돈 대신 맡겨 두는 재산)를 줄이지 않으면서 자기 돈을 더 많이 쓰이게 하고 싶어 해요 \parencite{bartoletti2021lending}. 그리고 2020년의 위기는 위험 설정값을 늦게 고치면 얼마나 큰 비용이 드는지 보여 줬어요 \parencite{gudgeon2020crisis}. 거래 상대가 계약이라면, 이런 서비스가 걸러 내며 묻는 질문은 그 계약이 강제로 청산될지예요. 그 결과 값에서, 허락한 주인 수가 송금 수보다 앞서는 것은 주된 기간에서는 가려졌지만 W0에서는 가려지지 않았어요. 페이지랭크 짝은 가려지지 않았어요. 그리고 이 신호는 거래를 잘하는 데에서는 더 낮고, 기간마다 계약 트레이더 몇십 개에만 기대고 있어요(\ref{sub:finding-neutral}절). 이 신호에 바탕을 둔 걸러 내기를 하려면, 더 큰 표본에서 등록한 시험이 필요해요. 이 점수들은 채무 불이행을 예측하는 도구가 아니에요.
\item 지갑 프로그램과 보안 도구. 사용자와 통합 개발자(여러 서비스를 이어 붙여 쓰는 개발자)는 스펜더 계약이 널리 허락받았는지 알고 싶어 해요. 스펜더를 허락한 주인이 몇이나 되는지, 또는 그 스펜더의 엔도스랭크를 보여 주는 배지(표시)는 그 위임을 분명히 드러내 줘요 \parencite{vogelsteller2015erc20,lundfall2020permit}. 이 배지는 그래프에서 점의 자리를 알려 줄 뿐이에요. 보안 감사도 아니고, 나쁜 계약을 막아 준다는 보장도 아니에요.
\item 유동성 공급자와 마켓 메이커. 이들은 가진 물량과 사고파는 값의 차이(스프레드)를 다스리려고, 송금 중심지를 찾고 흐름이 어떻게 짜여 있는지 봐야 해요. 그 목적에는 송금 수와 AWP가 알맞은 도구예요 \parencite{do2023awp}. 그 이로움은 걸러 내는 비용이 낮아지는 것이지, 손익이 좋아지는 것이 아니에요.
\item 에어드롭(공짜 토큰 나눠 주기) 프로그램과 거버넌스(운영 방향을 투표로 정하기) 프로그램. 시빌 무리는 토큰 나눠 주기와 거버넌스 참여 자격을 늘 악용해요 \parencite{douceur2002sybil,luo2025airdrop}. 여기서 정의한 엔도스랭크는 순위의 꼭대기 아래에서는 보호를 거의 해 주지 못해요. 닫힌 농장은 자기 주소 수에 맞는 몫의 약 세 배를 모으고, 그것을 한 목표에 몰아줄 수 있기 때문이에요. 그래서 주소 약 25개면 한 계약을 코호트의 상위 10분의 1 안으로 올릴 수 있어요. 다만 꼭대기에 오르려면 주소가 약 47{,}000개 필요해요(\ref{sec:sybil-cost-model}절). 이 수들 뒤에 있는 한계는 C-PR에도 그대로 들어맞아요. 프로그램은 서로 따로 활동해 온 주인들에게서 다양한 보증을 받으라고 요구할 수 있어요. 닫힌 무리 안에서만 도는 몫을 깎을 수도 있어요. 또는 순간 이동 몫과 매달린 점의 몫을 미리 확인한 씨앗 집합에만 보낼 수도 있어요. 하지만 이 가운데 어느 것도 정직한 주인에게서 보증을 사는 공격자는 막지 못해요.
\item 분석 회사와 규정 준수 회사. 거래소, 펀드, 규정 준수 팀은 바깥 사람도 다시 계산할 수 있는 거래 상대 순위가 필요해요. C-PR, 엔도스랭크, AWP는 정해 둔 기간과 설정값으로 공개하고 감사받을 수 있어요. 이 점이 회사만 아는 기계 학습 모형과 달라요 \parencite{packin2024credit}. 다만 시빌 공격에 얼마나 드러나 있는지를 밝히고, 이익 예측으로 내세우지 않아야 해요.
\end{itemize}

이 장면들 가운데 어디에서도 평판이 담보를 대신하지 않아요.

\dissertationSubsubheading[sub:deploy-multiples]{함께 쓰는 기반으로서의 표준 점수}

여러 쪽이 함께 쓰는 점수는, 그 하나하나가 점수를 다시 계산할 수 있고 점수가 어떻게 조작될 수 있는지 볼 수 있을 때만 쓸모가 있어요. C-PR, 엔도스랭크, AWP는 공개된 로그와 공개된 하루 가격 기록으로, 공개된 처리 과정을 써서 다시 계산할 수 있어요(부록~\ref{ch:appendix-eval}). 그리고 \ref{sec:sybil-cost-model}절과 \ref{sec:regulatory-ethics}절은 그 점수들이 어떻게 조작될 수 있는지 적어 두었어요.

\dissertationSubheading[sec:societal-impact]{사회와 실제에 주는 영향}

디파이(DeFi, 탈중앙 금융) 시장에서는 거래 상대를 확인하는 데 돈이 많이 들어요. 그리고 그것을 할 수 있는 힘은 고르게 나뉘어 있지 않아요. 돈이 많은 쪽은 회사만 아는 위험 도구에 돈을 낼 수 있을지 몰라요. 반면에 개인 참여자는 떠도는 이야기에 기대거나, 확인을 아예 건너뛰어요. 사용자들이 이미 해 준 허락으로 만들고, 다시 계산하는 비용도 싼 점수가 있다고 해요(표~\ref{tab:benchmark-runtime}). 이런 점수는 그 허락을 요청하는 계약을 확인하는 문턱을 낮춰 줘요. 그리고 어떤 계약이 받은 허락을 나쁘게 쓰면 주인이 몇이나 피해에 드러날지 보여 줘요. 송금 기록은 이것을 보여 주지 못해요. 예전에 해 준 허락은 토큰이 털려 나가는 길 가운데 하나예요(\ref{sec:approval-security}절). 하지만 이 점수가 털어 가기를 미리 내다본다는 것은 보이지 않았어요. 털어 가기 어림 규칙(drain heuristic, 토큰을 빼 가는 일을 대강 찾아내는 규칙)의 값은, 송금에 바탕을 둔 점수까지 포함해 모든 점수와 함께 올라갔어요(\ref{sec:holdout-results}절).

결과는 점수 하나를 권하기보다, 결정마다 알맞은 점수를 맞추는 쪽을 뒷받침해요. 그렇게 하면 흐름 점수로 허락을 정하는 위험도 줄어요. 위에서 인용한 차수 대비들은 라벨을 본 뒤에 계산했어요.

\textcite{packin2024credit}에서 다룬 탈중앙 신용 점수 시스템들은 속이 들여다보이지 않아요. 이와 달리 공개된 이벤트와 공개된 가격으로, 버전을 매긴 설정값을 써서 만든 점수는 설명하는 정보를 줘요. 규제 기관은 디파이 서비스를 분석할 때 그런 정보를 요구해야 한다고 제안해요 \parencite{zetzsche2020defi}. 이런 점수에는 의무도 따라와요. 순서만 따지고 그래프에 바탕을 둔 통계가 신원, 신용 기록, 또는 보안 감사로 잘못 받아들여지지 않게 해야 해요. 이 의무는 \ref{sec:regulatory-ethics}절에서 설명해요.

가치를 어떻게 나눌지에 영향을 주는 평판 지표는 무엇이든 악용될 수 있어요. 다음 절은 허락 화살표 위를 걷는 걷기를 조작하는 데 얼마나 드는지 모형으로 보여 줘요. 엔도스랭크와 C-PR도 그런 걷기예요.

\dissertationSubheading[sec:sybil-cost-model]{허락 그래프에 대한 시빌 공격 비용 모형}

시빌 공격은 값싼 가명을 써서, 믿을 만하다는 거짓 인상을 만들어요 \parencite{douceur2002sybil}. 송금 그래프에서는 나쁜 뜻을 가진 사람이 자기가 움직이는 계정들 사이에서 자산을 옮겨, 입차수와 들어오는 흐름을 억지로 부풀릴 수 있어요. 그 대가로 거래 수수료를 내고, 가진 물량을 잃을 수도 있어요 \parencite{cong2023washtrading}. 허락 그래프는 공격받는 곳(공격 표면)이 달라요. 각 화살표가 실제 송금이 아니라, 권한을 준 일을 기록하기 때문이에요. 이 절에서는 가짜 허락으로 허락 화살표 위의 걷기를 부풀리는 데 드는 비용을 간단한 모형으로 세워요. 계산은 엔도스랭크의 그래프에서 해요. 이 모형이 주는 한계는 C-PR에도 그대로 들어맞아요(\ref{sub:sybil-bound}절). 이 모형은 \textcite{gyongyi2005alliances}의 최적 링크 농장(점수를 가장 크게 부풀리도록 짠 링크 무리)을 바탕으로 해요. 그리고 \textcite{cheng2005sybilproof}의 결과도 바탕으로 해요. 그 결과에 따르면, 페이지랭크 같은 대칭 함수(점의 이름을 바꿔도 같은 규칙으로 계산하는 함수)는 시빌 공격을 막아 내지 못해요(\ref{sec:reputation-systems}절). 이 모형은 여기에 허락 화살표 하나의 비용과, 이 연구의 그래프에서 얻는 이득의 크기를 보태요.

\dissertationSubsubheading[sub:sybil-setup]{공격 설정}

공격자가 주소 집합 $\mathcal{S}$를 움직인다고 해요. 그리고 보증 그래프 $G_{\text{approve}}=(V,E,w)$에서 정직한 주소들을 $\mathcal{H}=V\setminus\mathcal{S}$라고 써요. 무게가 $w_{uv}>0$인 화살표 $(u,v)$는 $u$가 자기 토큰을 쓸 권한을 $v$에게 주었다는 것을 기록해요. 이 무게는 식~\eqref{eq:approve-weight}의 무게예요. 감쇠 계수 $d\in(0,1)$를 쓰면, 엔도스랭크는 $G_{\text{approve}}$ 위의 페이지랭크 벡터 $\mathbf{r}$를 다음 식의 해(답)로 구해요 \parencite{page1999pagerank}.
\begin{equation}
\label{eq:pagerank-sybil}
\mathbf{r} = d\, P^{\top} \mathbf{r} + \frac{1-d}{|V|}\mathbf{1},
\end{equation}
여기서 $P$는 식~\eqref{eq:transition}의 이동 행렬이에요. 나가는 화살표가 있는 점은 자기 행(가로줄)을 나가는 허락 무게에 비례해서 나눠요. 매달린 점은 나가는 화살표가 없는 점이에요. 매달린 점의 행은 모든 칸이 $1/|V|$예요. 그래서 그 몫은 모든 점에 고르게 퍼져요 \parencite{langville2006pagerank}. 이 고정점은 식~\eqref{eq:pagerank}의 반복이 닿는 고정점과 같아요. 구현한 코드는 그 반복을 돌려요. 그리고 그 반복은 매달린 점의 몫을 고른 순간 이동 벡터를 통해 다시 넣어 줘요. 공격자는 목표 집합 $\mathcal{G}\subseteq\mathcal{S}$의 평균 엔도스랭크를 가장 크게 만들려고 해요.
\begin{equation}
\label{eq:objective}
J(\mathcal{G}) = \frac{1}{|\mathcal{G}|}\sum_{i\in\mathcal{G}} r_i,
\end{equation}
다만 공격 비용에는 예산 제약(쓸 수 있는 돈의 한도)이 있어요. 공격 비용은 바로 다음에 정의해요.

\dissertationSubsubheading[sub:sybil-cost]{비용의 구성 요소}

공격자는 화살표 집합 $E_{\mathcal{S}}$를 만들어요. 이 화살표들은 모두 $\mathcal{S}$에서 출발해요. 공격자는 정직한 주인에게서 $\mathcal{S}$로 들어오는 화살표 집합 $E_{\mathcal{H}\to\mathcal{S}}$를 살 수도 있어요. 비용은 점수 기간 $q$개에 걸쳐 세요. 이 기간 수는 점수를 새로 계산하는 횟수이고, 목표들은 그동안 순위를 지키려고 해요. 모든 화살표에 빠짐없이 드는 비용은 가스뿐이에요.

\begin{enumerate}
\item 가스 비용. \texttt{approve}나 \texttt{permit}을 부를 때마다 가스가 들어요. 허락을 한 번 새로 고치는 데 드는 수수료의 기댓값(평균적으로 예상되는 값)을 $c_{\mathrm{gas}}$라고 써요. 엔도스랭크는 각 기준 시점(점수를 고정하는 때) 전 열두 달 동안의 \texttt{Approval} 이벤트를 읽어요. 그래서 어떤 허락은 그 허락이 기간 안에 드는 그래프에만 들어가요. 닫힌 농장 안에서는 허락의 시간 무게가 중요하지 않아요. 행 정규화(각 행의 합이 1이 되게 나누기)가 그것을 없애기 때문이에요(항목 2). 그래서 닫힌 농장은 적어도 열두 달에 한 번은 허락을 다시 기록해야 해요. 그리고 점수를 새로 계산하는 일이 그런 기간 $y$개에 걸치면, $C_{\mathrm{gas}} = y\,|E_{\mathcal{S}}|\,c_{\mathrm{gas}}$예요.

\item 노출 비용(토큰을 남이 쓸 수 있게 열어 두어서 생기는 비용). 허락은 토큰을 하나도 묶어 두지 않아요. 그리고 주인은 허락할 때 그 토큰을 조금도 가지고 있지 않아도 돼요. 금액을 크게 적어도 돈이 들지 않아요. 무제한 허락(금액 한도가 없는 허락)은 뒤에 잔고가 하나도 없어도 금액 무게 $V$(금액이 아무리 커도 1을 넘지 않게 바꾸는 식)를 가득 받아요. 그리고 행 정규화를 하고 나면, 나가는 화살표가 하나뿐일 때는 그 무게가 얼마든 $P_{uv}=1$이 돼요. 다섯 토큰의 목록도 이 비용을 올리지 않아요. 목록에 있는 토큰을 허락하는 데에도 잔고가 필요 없기 때문이에요. 그래서 $\mathcal{S}$ 안의 화살표는 아무것도 드러내 놓지 않고, 가스 말고는 비용이 들지 않아요. 노출은 토큰을 가진 공격자 주소가 스펜더 $v\notin\mathcal{S}$를 허락할 때만 생겨요. 그러면 $v$가 그 토큰을 옮길 수 있기 때문이에요. 공격자는 예를 들어 잘 알려진 계약을 쓰는 평범한 사용자처럼 보이려고 그렇게 할 수 있어요. 이런 화살표는 점수의 몫이 $\mathcal{S}$ 밖으로 빠져나가게도 해요. 이런 노출 하나가 한 기간 동안 공격자에게 주는 비용의 기댓값을 $\ell_{uv}$라고 하면, $C_{\mathrm{exp}} = q\sum_{(u,v)\in E_{\mathcal{S}},\,v\notin\mathcal{S}} \ell_{uv}$예요. 닫힌 농장에서는 이 값이 0이에요.

\item 허락 취소와 들킴 비용. 정상적인 사용자는 오래된 허락을 가끔 거둬요. 그리고 서로만 허락해 주는 새 주소들의 농장은 들킬 수도 있어요. 들키는 비용과, 들키지 않으려고 공격자가 하는 허락 취소의 비용을 합쳐, 기간 하나당 그 기댓값을 $c_{\mathrm{rev}}$라고 해요. $q$개 기간 동안에는 $C_{\mathrm{rev}} = q\,c_{\mathrm{rev}}$예요.

\item 사들인 정직한 화살표. 공격자는 정직한 주인에게 돈을 주고, 자기 주소들을 허락하게 할 수도 있어요. 예를 들어 사용자에게 권한을 달라고 하는 앱을 쓸 수 있어요. $h\in\mathcal{H}$이고 $a\in\mathcal{S}$인 사들인 화살표 $(h,a)$는, $h$가 마지막으로 허락한 뒤 열두 달 동안, 그것도 그 허락이 살아 있는 동안에만 그래프에 남아요. 그리고 시간이 지날수록 $h$의 행에서 차지하는 몫이 줄어요. 그래서 그 화살표가 계속 효과를 내게 하려면, 기간 하나당 다음 비용이 들어요.
\begin{equation}
\label{eq:marginal}
c_{\mathrm{in}}(h,a) = c_{\mathrm{gas}} + p_h + \rho_h,
\end{equation}
여기서 $p_h$는 정직한 주인에게 주는 돈이에요. 그리고 $\rho_h$는 위험 보상금(위험을 떠안는 대가로 더 받는 돈)이에요. 주인은 공격자가 그 허락을 써 버릴 위험에 대비해서 이 돈을 요구해요.
\end{enumerate}

$q$개 기간 동안 공격의 전체 비용은 다음과 같아요.
\begin{equation}
\label{eq:total-cost}
C = C_{\mathrm{gas}} + C_{\mathrm{exp}} + C_{\mathrm{rev}} + q\sum_{(h,a)\in E_{\mathcal{H}\to\mathcal{S}}} c_{\mathrm{in}}(h,a).
\end{equation}
닫힌 농장에서는 $C_{\mathrm{exp}}=0$이고 $E_{\mathcal{H}\to\mathcal{S}}$는 비어 있어요. 그래서 $C=y\,|E_{\mathcal{S}}|\,c_{\mathrm{gas}}+q\,c_{\mathrm{rev}}$예요.

\dissertationSubsubheading[sub:sybil-bound]{영향력과 비용 효과의 한계}

페이지랭크에 대한 흔한 민감도 결과(조건을 바꾸면 값이 얼마나 움직이는지에 관한 결과)에 따르면, 화살표 하나는 $\mathbf{r}$를 움직여요. 그 크기는 출발점의 정상 분포 몫에 비례하고, 그 화살표의 정규화한 나가는 무게에도 비례해요 \parencite{langville2006pagerank}. 먼저 닫힌 농장을 생각해 봐요. 닫힌 농장에서는 $\mathcal{S}$의 모든 점이 나가는 화살표를 적어도 하나 가져요. 이 화살표들은 모두 $\mathcal{S}$ 안에서 끝나요. 그리고 어떤 정직한 주인도 $\mathcal{S}$ 안의 주소를 허락하지 않아요. 매달린 점의 집합을 $D$라고 해요. 그러면 매달린 점은 모두 $\mathcal{H}$ 안에 있어요. 그리고 그 점들의 정상 분포 몫을 $r_D=\sum_{j\in D} r_j$라고 해요. 식~\eqref{eq:pagerank-sybil}에 따르면, 점 $i\in\mathcal{S}$는 농장에서 $d\sum_{j\in\mathcal{S}} r_j P_{ji}$를 받아요. 순간 이동에서는 $(1-d)/|V|$를 받고, 매달린 점의 행에서는 $d\,r_D/|V|$를 받아요. 농장 점들의 행은 $\mathcal{S}$ 안에서 더하면 1이 돼요. 그래서 $\mathcal{S}$ 전체에 대해 더하면 $\sum_{i\in\mathcal{S}} r_i = d\sum_{i\in\mathcal{S}} r_i + |\mathcal{S}|\,(1-d+d\,r_D)/|V|$가 돼요. 이것을 정리하면 다음과 같아요.
\begin{equation}
\label{eq:closed-cluster}
\sum_{i\in\mathcal{S}} r_i = \frac{|\mathcal{S}|}{|V|}\cdot\frac{1-d+d\,r_D}{1-d}.
\end{equation}
첫째 인수(곱해지는 값)는 농장의 주소 수 몫(전체 점 가운데 농장 주소가 차지하는 비율)이에요. 둘째 인수는 농장이 그 몫의 몇 배를 모으는지 말해 줘요. 둘째 인수는 $r_D=0$일 때만 1이에요. 곧 매달린 점이 없는 그래프에서만 그래요. 그렇지 않으면 매달린 점의 행이 정직한 쪽의 몫 일부를 농장에 넘겨줘요. 매달린 점은 $\mathcal{S}$ 밖에 있으므로 $r_D\le 1-\sum_{i\in\mathcal{S}} r_i$예요. 이것을 식~\eqref{eq:closed-cluster}에 넣으면, 어떤 그래프에서든 성립하는 한계(넘을 수 없는 값)를 얻어요.
\begin{equation}
\label{eq:cluster-mass}
\sum_{i\in\mathcal{S}} r_i \le \frac{|\mathcal{S}|}{(1-d)\,|V| + d\,|\mathcal{S}|}.
\end{equation}
$|V|$에 비해 작은 농장이라면, 이 한계는 주소 수 몫의 $1/(1-d)$배에 가까워요. $T_{\mathrm{obs}}$의 매칭 코호트 엔도스랭크 그래프에는 점이 $2{,}513{,}982$개 있어요. 이 그래프에서 처리 과정에 들어 있는 페이지랭크 코드로 계산하면, $d=0.75$, $0.85$, $0.95$일 때 $r_D=0.365$, $0.394$, $0.419$예요. 그래서 이 그래프의 닫힌 농장은 자기 주소 수 몫의 2.10배, 3.23배, 8.96배를 모아요. 같은 코드로 직접 점검해도 결과가 같아요. 주소 20개를 고리(둥글게 이은 모양)로 이어서, 각 주소가 다음 주소를 허락하게 하고, 그래프에 더해 봤어요. 그러면 이 고리는 $d=0.85$에서 자기 주소 수 몫의 3.23배를 모아요.

식~\eqref{eq:cluster-mass}은 농장 전체의 몫을 제한할 뿐, 목표 하나의 점수를 제한하지는 않아요. 그리고 공격자는 그 몫을 한곳에 몰 수 있어요. 별 모양을 생각해 봐요. 농장 주소 $f$개가 저마다 목표 $t$를 허락하고, $t$도 그 주소 하나하나를 다시 허락해요. 모든 농장 주소 $a$에게는 $t$가 유일한 나가는 이웃이에요. 그래서 $P_{at}=1$이에요. 그리고 $t$의 행은 농장 위에서 더하면 1이 돼요. 각 점이 순간 이동과 매달린 점의 행에서 받는 몫을 $c=(1-d+d\,r_D)/|V|$라고 써요. 그러면 균형 식(점이 받는 몫을 모두 더한 식) $r_a=c+d\,r_t P_{ta}$와 $r_t=c+d\sum_a r_a$에서 다음 식을 얻어요.
\begin{equation*}
r_t = \frac{c\,(1+df)}{1-d^2}.
\end{equation*}
$d=0.85$에서는 $c=1.93\times10^{-7}$이에요. 이것은 허락을 하나도 받지 않은 모든 주소가 함께 가지는 점수예요. $f=10$이면(주소 열한 개, 허락 스무 개), 별 모양을 더한 그래프에서 목표는 $r_t=6.60\times10^{-6}$에 이르러요. 이 값은 코호트 계약들의 엔도스랭크 중앙값(한가운데 값)보다 높아요. 하지만 가장 높은 점수($0.0278$)의 $2.4\times10^{-4}$에 지나지 않아요. 점수는 $f$에 따라 일정한 비율로(직선으로) 커져요. 농장 주소를 하나 더할 때마다 $r_t$에 $d\,c/(1-d^2)$가 더해지고, 허락이 두 개 더 들어요. 이 식을 $f$에 대해 풀면, 코호트의 90번째 백분위(작은 것부터 줄 세웠을 때 90\% 자리의 값)에 이르는 데 농장 주소가 약 25개 필요해요. 99번째 백분위에는 약 2{,}292개, 코호트에서 가장 높은 점수에는 약 47{,}073개가 필요해요. 그러니까 평범해 보이는 데는 허락 몇십 개면 되지만, 꼭대기에 이르려면 허락이 약 94{,}000개(농장 주소 하나에 두 개씩) 필요해요. 순위가 가장 높은 코호트 계약이 그래프 전체 정상 분포 몫의 약 3\%를 가지고 있기 때문이에요. 그래프는 각 동결일 전 열두 달 동안의 허락만 담아요. 그래서 농장은 목표의 순위를 지키려면 적어도 한 해에 한 번은 이 허락들을 다시 기록해야 해요(\ref{sub:sybil-cost}절). 이 연구는 $c_{\mathrm{gas}}$를 재지 않았어요. 그래서 이 수들을 값(돈)으로 바꾸지 않아요.

그러니까 순위의 윗부분에 오르는 데 정직한 들어오는 화살표는 필요 없어요. 하지만 사들인 정직한 화살표는 이득을 더 키워요. 농장에서 정직한 주소로 나가는 화살표가 없다고 해요. 그러면 정직한 주인 $h$에게서 사들인 화살표 $(h,a)$는 농장의 몫에 대략 다음만큼을 더해요(1차 근사, 곧 아주 작은 변화만 따진 값).
\begin{equation*}
\Delta\sum_{i\in\mathcal{S}} r_i \approx \frac{d\,r_h\,P_{ha}}{1-d}
\end{equation*}
여기서 $P_{ha}$는 새 화살표가 $h$의 행에서 차지하는 몫이에요. $h$에게서 새로 사들인 크거나 무제한인 허락은 $V(z)$가 1에 가까워서, $h$의 행에서 약 $\sigma(0)/(o(h)+\sigma(0))$만큼을 차지해요. 그래프에서 중앙값에 있는 주인에게 그 몫은 $0.683$이에요. 그래서 공격자가 사려는 화살표는, 자기 허락이 오래되었거나 작거나 적은 주인에게서 받는 크거나 무제한인 허락이에요. 될 수 있으면 $r_h$가 높은 주인이 좋아요. 이런 허락을 해 주면, 주인이 가진 그 토큰 잔고를 공격자가 마음대로 쓸 수 있게 돼요. 그래서 식~\eqref{eq:marginal}의 위험 보상금 $\rho_h$는 주인이 가진 양이 많을수록 커지고, 그 토큰을 하나도 가지지 않은 주인에게서만 0에 가까워요. 상한을 둔 C-PR의 허락 층에서는, 이런 허락이 주인의 다른 허락들에 비해 클 때만 행에서 큰 몫을 차지해요(\ref{sub:coupled-graph}절). 공격의 비용 효과(비용 한 단위로 얻는 이득)는 다음과 같아요.
\begin{equation}
\label{eq:efficiency}
\eta = \frac{\Delta J(\mathcal{G})}{C},
\end{equation}
여기서 $\Delta J(\mathcal{G})$는 공격 화살표들이 만들어 내는, 식~\eqref{eq:objective}의 목표 값의 증가분이에요. 이 증가분은 식~\eqref{eq:total-cost}의 $C$가 값을 치르는 $q$개 기간 동안 유지돼요.

식~\eqref{eq:closed-cluster}--\eqref{eq:efficiency}을 보면, 여기서 정의한 엔도스랭크는 닫힌 농장 때문에 싸게 부풀려질 수 있어요. 고른 순간 이동과, 매달린 점의 몫을 고르게 다시 나누는 방식 아래에서는, 닫힌 무리가 자기 주소 수 몫에 묶여 있지 않아요. 결합 걷기도 같은 농장에 똑같이 드러나 있어요. 식~\eqref{eq:cluster-mass}을 끌어낼 때는 고른 순간 이동과 매달린 점의 몫을 고르게 퍼뜨리는 것만 써요. 그리고 C-PR도 이 두 가지를 똑같이 써요. 나가는 화살표가 허락뿐인 농장 주소는, $\lambda$가 얼마든 그 허락 화살표를 무게 1로 따라가요(식~\eqref{eq:coupled-transition}). 그리고 나가는 화살표가 하나뿐이면, 그 금액이 얼마든 그 화살표가 행 전체를 차지해요. 그래서 상한 $U$도 C-PR을 지켜 주지 못해요. 위의 농장 배수(농장이 주소 수 몫의 몇 배를 모으는지)와 별 모양 크기는 허락 그래프에서만 계산했어요. C-PR의 합친 그래프에서는 매달린 점의 몫이 다르고, 그래서 $c$도 달라요. 그 값들은 계산하지 않았어요. 닫힌 무리를 정말로 묶어 두는 방법은 걷기 자체를 바꿔요. 트러스트랭크처럼, 순간 이동과 매달린 점의 몫을 미리 확인한 씨앗 집합에만 보낸다고 해요 \parencite{gyongyi2004trustrank}. 그러면 그 집합 밖의 무리는 정직한 주인에게서 오는 화살표를 통해서만 몫을 얻어요. 그리고 공격자는 그런 화살표를 식~\eqref{eq:marginal}의 값을 치르고 사야 해요. AWP의 다시 출발 규칙도 송금 쪽을 지켜 주지 못해요. 주소들끼리 적은 금액을 주고받는 농장은 그 주소들의 활발함(활동 정도)을 높여요. 그래서 다시 출발 몫을 자기 쪽으로 끌어와요. 활동 재시작을 쓴 엔도스랭크도 같은 약점을 물려받아요. 시빌 보정 안정성 대리 지표(\ref{sub:proxy-sybil}절)는 몇몇 주소에서 되풀이되는 송금은 잡아내지만, 새로 만든 주소들의 농장은 잡아내지 못해요. 그래서 이 지표는 적대적 평가(공격자 입장에서 실제로 공격해 보는 평가)를 대신할 수 없어요. 적대적 평가는 앞으로의 연구로 남겨 둬요.

\dissertationSubheading[sec:regulatory-ethics]{규제와 윤리에 대한 고려}

온체인 평판 지표는 투명성(속을 들여다볼 수 있는 정도)에 대한 질문과 형평성(모두에게 공평한지)에 대한 질문을 함께 던져요. 앞선 법학 연구에 따르면, 탈중앙 금융 시장은 결제가 온체인에서 이루어져도 여전히 공시(정보를 공개해 알리는 일) 규칙과 소비자 보호 규칙에 묶여 있어요 \parencite{zetzsche2020defi}. \textcite{packin2024credit}도 나름의 경고를 보태요. 모형을 해석할 수 없거나 모형에 이의를 제기할 수 없을 때, 탈중앙 신용 점수는 블랙박스(속을 들여다볼 수 없는 장치)의 해를 계속 이어 갈 수 있어요. 모형이 정당한 절차 없이 신원 데이터베이스와 이어질 때도 그래요. C-PR, 엔도스랭크, AWP는 모두 온체인 이벤트와 공개된 하루 가격 기록으로, 공개된 알고리즘을 써서 계산해요. 계산은 늘 같은 결과가 나오는 방식(결정론적)이에요. 그래서 누구나 공개된 소프트웨어(오픈 소스)로 다시 계산할 수 있어요(부록~\ref{ch:appendix-eval}). 이런 공개성이 블랙박스라는 비판에 어느 정도 답해 줘요. 하지만 누가 이익이나 손해를 보는지에 대한 걱정(분배 문제)은 그대로 남아요. 활동이 적은 계약은 코드가 어떻든 점수가 더 낮게 나올 것이기 때문이에요. 새 계약도 여기에 들어요.

윤리적으로 보면, 이 점수들을 보안 등급으로 쓰면 절대 안 돼요. 전통적인 뜻의 신용도(돈을 갚을 만한 정도)를 재는 값으로 써서도 안 돼요. 그리고 한 계약을 허락한 주인들이 서로 다른 사람이라는 증거로 써서도 안 돼요 \parencite{douceur2002sybil}. 이 점수들은 정의한 그래프 짜임(위상) 안에서 계약이 어디에 있는지를 순위로 나타내는 지표예요. 실제로 쓰는 곳에서 점수 기준값에 따라 금융 서비스를 내준다면, 점수 뒤에 있는 그래프 정의, 기간, 한계를 분명히 문서로 밝혀야 해요. 무엇보다 시빌 공격에 드러나 있다는 점을 밝혀야 해요. 해석할 수 있는지(해석 가능성)를 가운데에 둔 규제 틀이라면, \textcite{kotb2023credit}에서 살펴본 것 같은, 회사만 아는 기계 학습 모형보다 이런 순위를 더 좋아할 거예요. 다만 점수를 내놓는 쪽이 대출 상환율을 얼마나 잘 예측할 수 있는지 부풀려 말하지 않아야 해요 \parencite{packin2024credit}.

열린 블록체인에서 개인 정보 보호(프라이버시)는 장부의 성질 자체 때문에 한계가 있어요 \parencite{nakamoto2008bitcoin}. 그리고 평판 점수를 계산하는 데는 공개된 자료 말고는 아무것도 쓰지 않아요. 각 주소 뒤에 있는 신원을 더 캐내려는 시도는 하지 않기를 권해요. 그리고 점수를 지문처럼 영원히 따라다니는 식별자(누구인지 알아보는 표시)로 다뤄서는 안 돼요. 이 연구는 계약 주소에만 점수를 매기고, 어떤 주소도 실제 세상의 개인과 이으려고 하지 않아요(\ref{sec:ethics}절).

\dissertationSubheading[sec:threats-validity]{타당성을 위협하는 것들}

여기서 보고한 발견은 아래의 한계에 비추어 읽어야 해요. 이 한계들은 구성 타당도, 내적 타당도(비교가 정말 그 차이만 가려내는지), 외적 타당도(바깥 자료로도 맞는지), 적대적 타당도(누가 일부러 속이려 해도 맞는지)로 나누어 묶었어요.

\dissertationSubsubheading[sub:threat-construct]{구성 타당도}

아래의 위협들은, 점수와 점검이 잰다고 읽히는 것을 정말 재는지에 관한 거예요.

\begin{itemize}
\item 이 순위 점수는 신용 등급이 아니에요. 담보 없는 대출에서 나온 채무 불이행 라벨은 분석에 하나도 들어가지 않아요.
\item 같은 기간 점검은 가까운 주변(국소)의 허락 통계와 송금 통계에 바탕을 둬요.
\item 홀드아웃 구성개념. 나중의 새 허락도 여전히 보증 이벤트이고, 나중의 새 송금자도 흐름 이벤트예요. 둘 다 점수 기간 밖에 있지만, 어느 쪽도 신용이나 거래를 따로 재는 구성개념은 아니에요. 그리고 둘 다 새로 들어온 것을 센 수라서, 지금의 거래 상대 수를 따라가는 점수라면 어느 것에든 유리해요. 두 기간 모두에서 이 라벨들을 가장 잘 줄 세우는 것도 그냥 센 수예요(\ref{sub:finding-rq3}절).
\item 분석 단위. 허락 위를 걷는 걷기들은 허락을 받는 쪽에 점수를 매겨요. 그리고 여기서 점수를 매긴 주소는 모두 계약이에요(\ref{sub:cohorts}절). 트레이더의 평판에 대한 주장은 W0 트레이더 분석과 \ref{sub:finding-neutral}절의 비교에 기대요. 둘 다 작은 표본에서 한 거예요.
\item 자기 구성개념과의 일치는 어느 정도 기계적으로 생겨요. 엔도스랭크는 스펜더의 받은 허락 수를 그래프 전체에 걸쳐 매끄럽게 다듬어요. 그리고 AWP는 금액으로 무게를 단 들어오는 송금에 같은 일을 해요. 그래서 각 점수는 만든 방식 때문에 저절로 자기 가족과 맞을 것으로 예상돼요. 표~\ref{tab:tau-diff}에서 한 점수 안의 가족끼리 견준 대비는 이 까닭으로 계산했어요. 동점이 적어서 이 대비는 같은 조건끼리 견줘요. 하지만 이 대비는 여전히 도구에 대한 점검이에요.
\item 금액 무게. 엔도스랭크와 AWP는 각 달러 금액을 $V(z)$에 넣어요. 이때 기울기는 중앙값만큼의 금액이 무게 절반을 갖도록 정해요(\ref{sec:amount-revision}절). 그래서 무게는 화살표가 담은 허락과 송금의 개수와 나이뿐 아니라, 그 금액도 기록해요. 기울기는 수정에서 고른 값이고, 같은 기간에서만 점검했어요(표~\ref{tab:robustness-slope}). C-PR의 층들은 $V$의 한계 없이 달러 금액을 더하고, 각 허락은 $U$에서 잘라요(\ref{sub:coupled-graph}절). 두 스펜더 홀드아웃 모두에서, 엔도스랭크는 두 라벨을 허락 층만 걷는 것보다 조금 더 잘 줄 세워요(표~\ref{tab:holdout-spenders}와~\ref{tab:fresh-posthoc}). 짝지은 대비는 계산하지 않았어요. 그래서 금액에 무게를 다는 방식이 이 구성개념에 중요할 수도 있어요. 그리고 어느 무게 다는 방식도 옳다고 밝혀지지 않았어요. 이런 무게 선택은 $\mathbf{r}$와 순위 상관으로 퍼져 들어가요 \parencite{langville2006pagerank}.
\end{itemize}

\dissertationSubsubheading[sub:threat-internal]{내적 타당도}

아래의 위협들은, 비교의 결과를 만들어 냈을 수도 있는 선택에서 비교가 자유로운지에 관한 거예요.

\begin{itemize}
\item 함께 쓰는 점수 기간. 모든 방법과 대리 지표는 동결일 전 열두 달 동안에 대해 계산해요. 같은 기간과 W1에서는 2025년 7월 1일부터 2026년 6월 30일까지이고, W0에서는 2025년 4월 1일부터 2026년 3월 31일까지예요. 그 달들에만 있었던 시장 분위기(값의 출렁임, 보상 프로그램, 서비스 변화)가 일치도를 올리거나 내릴 수 있어요. 기간이 시작되기 전에 마지막으로 정한 허락은, 아직 살아 있어도 화살표가 아니에요. 그래서 주인들이 오래전에 허락하고 다시 허락하지 않은 계약은 그 화살표들을 잃어요. 감쇠 계수 점검과 금액 기울기 점검은 이 설정값들에 대한 민감도를 재요. 하지만 그 가운데 어느 것도 달력 기간을 바꾸지 않아요.
\item 코호트 고르기. 매칭 코호트는 $T_{\mathrm{obs}}$ 전 열두 달 동안 다섯 토큰 가운데 하나의, 0이 아닌 허락을 적어도 세 주인에게서 받은 아비트럼 원의 계약 6,483개예요. 결과는 허락이 더 적은 계약이나, 다른 토큰의 허락만 받은 계약에는 그대로 옮겨지지 않을 수도 있어요. 그리고 그래프가 모든 코호트 계약을 보는 것도 아니에요(\ref{sub:cohorts}절). 표본 정의에 대한 민감도는 평가하는 무리가 바뀔 때 계수가 얼마나 크게 움직이는지 보여 줘요(표~\ref{tab:robustness-sample-size}).
\item 기준선 충실도(비교 기준을 원래 방법대로 따랐는지). AWP는 화살표 무게와 다시 출발에서 발표된 방법을 따라요(\ref{sub:awp-graph}절). 그 논문은 설정값을 알려 주지 않아요. 그래서 $k$, $t_0$, $b$는 이 연구가 고른 값이에요(위의 금액 무게 항목을 봐요). 그 논문은 이더리움 거래에 점수를 매겼어요. 여기서 AWP는 코호트의 에고 네트워크 위에서, 아비트럼 원의 ERC-20 토큰 다섯 개의 송금에 점수를 매겨요. 송금 금액은 달러로 매겨요. 에고 네트워크에서는 코호트 밖에 있는 보낸 쪽의 활발함이, 그 주소가 코호트 계약에 보낸 송금만으로 정해져요.
\item 분석 순서. 코호트, 점수, 라벨, 통계, 규칙 A는 2026년 10월 9일의 분석 계획에서 정했어요. 이 연구의 로그를 하나도 뽑기 전이었어요. 규칙 B와 그 허용 폭, 두 민감도 분석, 그리고 트레이더 결과에 대한 비교는 그날 늦게 등록했어요. W0 결과를 안 뒤였고, 이 연구가 7월부터 9월까지의 로그를 뽑기 전이었어요(\ref{sub:fresh-holdout}절). 다른 모든 비교는, 차수 대비들도 포함해서, 관련된 라벨을 본 뒤에 계산했고 구간과 함께 보고해요. 그리고 어느 것도 질문의 답을 정하지 않아요.
\item 규칙의 비교 대상. 규칙 A와 B는 C-PR을 자기 층을 하나만 걷는 것과 견줘요. 결합만 바꾸고 다른 것은 하나도 바꾸지 않는 대비는 이것뿐이에요. C-PR은 금액 무게와 고르게 다시 출발하기에서도 엔도스랭크나 AWP와 달라요. 그래서 C-PR을 그 점수들과 견준 대비(표~\ref{tab:cpr-contrasts})에는 결합의 효과와 그런 선택의 효과가 섞여 있어요. 발표된 걷기와 견주면서 규칙에 묶인 대비는, AWP를 비교 대상으로 삼은 등록한 민감도 분석 하나뿐이에요.
\item 의존성(서로 기대어 있음)과 누출(정보가 새어 들어옴). 주인이나 서비스를 함께 쓰는 계약들도 서로 관계없는 것처럼 다시 뽑아요. 그리고 사후 대비는 여러 번 비교한 것을 보정하지 않았어요(\ref{sec:statistical-tests}절). 동결일마다 코호트는 그 동결일 자신의 점수 기간의 허락에서 뽑아요. 그래서 어떤 라벨 기간도, 자기 자료를 뽑아 올 중심 계약들을 고르지 않아요(\ref{sec:holdout-protocol}절). 하지만 다섯 토큰은 $T_{\mathrm{obs}}$의 점수 기간 안에서 그 토큰들의 행 수(기록 수)에 따라 골랐어요. 그리고 그 점수 기간에는 W0의 라벨 기간이 들어 있어요.
\item 주 분석을 사후에 정한 것. 수정 분석(등록 뒤에 금액을 달러로 바꿔 다시 한 분석)의 금액, 토큰 목록, 점수 기간은 세 기간 모두에서 등록 분석의 결과를 안 뒤에 정했어요(\ref{sec:amount-revision}절). 그래서 이 논문의 주 분석인 수정 분석은 사후 분석이에요. 그리고 규칙 A, B, R1--R3에 따른 수정 분석의 판정은 등록한 시험이 아니라 재분석이에요. 수정 분석의 기울기와 상한은 W0 점수 기간의 분위수(자료를 크기순으로 나누는 자리의 값)에서 나와요. 이 점수 기간은 $t_1$에서 끝나요. 그리고 등록 분석도 규칙 A와 B, 그리고 R2와 R3에서 같은 판정에 이르렀어요. R1은 고친 점수에서는 성립하지만, 등록 분석에서는 성립하지 않았어요(부록~\ref{sec:map-registered}). 두 분석 사이에 다른 대비도 있어요. F3, 그리고 두 기간 모두에서 새 송금자에 대한 AWP 빼기 송금 입차수의 부호가 그래요. 이 값은 등록 분석에서는 0보다 위에 있었어요(표~\ref{tab:registered-vs-revised}). 그리고 W0에서 Q의 부호와, W0의 새 송금자에서 S-PR을 송금 걷기와 견준 값의 부호도 달라요.
\item 시간을 잰 환경. 모든 실행 시간은 컴퓨터 한 대에서 잰 실제 걸린 시간이에요(부록~\ref{ch:appendix-eval}). 표~\ref{tab:benchmark-runtime}의 표준편차는 한 설정을 되풀이해서 푼 것만 담아요. 그리고 감쇠 계수를 바꿔 가며 잰 실험에서는 그래프와 감쇠 계수의 짝마다 시간을 한 번만 쟀어요. 컴퓨터 사이의 차이는 재지 않았어요. 그래서 비교는 초 대신 비율로 적어요.
\end{itemize}

\dissertationSubsubheading[sub:threat-external]{외적 타당도}

아래의 위협들은, 결과를 살펴본 자료 너머로 얼마나 멀리 넓힐 수 있는지에 관한 거예요.

\begin{itemize}
\item 체인 하나와 기간 하나. 증거는 아비트럼 원 하나에서만 나와요. 다섯 토큰 가운데 하나로 적어도 세 주인이 허락한 계약들에서 나와요. 그리고 각 동결일 전 열두 달 동안 점수를 매긴 활동과, 석 달짜리 라벨 기간 두 개에서 나와요. 발견을 다른 롤업(여러 거래를 묶어 처리하는 체인), 레이어 1 체인(가장 바닥에 있는 블록체인), 대출 중심 서비스로 넓히려면, 여러 서비스에 걸친 검증이 필요해요. 디파이 서비스들은 서비스 설계가 서로 다르기 때문이에요 \parencite{werner2022sok}. 시장 구조도 서로 달라요 \parencite{schar2021defi}.
\item 표본 규모에 따른 변화와 일치도의 규모에 따른 변화. 부분 표본 단계들은 매칭 코호트 그래프의 부분 그래프에서 잰 실행 시간에만 관련이 있어요(\ref{sub:runtime-scaling}절). 다른 크기의 무리에서도 $\tau$의 모양이 그대로인지는 보여 주지 않아요.
\item 에고 네트워크의 깊이. 각 그래프는 코호트 계약에 닿는 화살표만 남겨요. 그래서 퍼뜨림에 대한 발견은 보증의 길이 더 긴 그래프로는 그대로 옮겨지지 않아요(\ref{sub:theory-closed-form}절).
\end{itemize}

\dissertationSubsubheading[sub:threat-adversarial]{적대적 타당도}

아래의 위협들은 공격자가 점수에 무엇을 할 수 있는지에 관한 거예요.

\begin{itemize}
\item 시빌 공격과 자전 거래. 공격자가 허락이나 송금의 모양을 꾸며서 순위를 조작할 수 있어요 \parencite{douceur2002sybil}. 시빌 보정 대리 지표는 이 문제를 일부만 다뤄요. 그리고 \ref{sec:sybil-cost-model}절의 비용 모형은 식으로 따진 분석이에요. 실제 자료로 한 레드팀 평가(일부러 공격하는 역할을 맡아 시험하는 평가)가 아니에요.
\item 전략적인 허락 행동. 허락으로 만든 점수를 실제로 쓰기 시작하면, 이익을 따지는 참여자들이 순위를 긁어모으려고 허락을 전략적으로 해 줄 수 있어요. 점수를 쓰기 전(배포 전)에 관찰 자료로 어림한 $\tau$는, 순위가 그런 행동을 얼마나 잘 버티는지를 부풀려 보여 줄 수 있어요.
\end{itemize}

\dissertationSubheading[sec:limitations]{한계}

이 표본에서 C-PR의 강점은 이래요. 두 기간 모두에서 두 스펜더 라벨을 엔도스랭크보다 위에 줄 세우는 것, 두 가족과의 일치가 자기 두 층 사이에 놓이는 것, 그리고 송금 걷기가 동점으로 남겨 두는 계약들 사이에 순서를 주는 것이에요. 약점은 이래요. 두 규칙을 모두 통과하지 못했어요. 나중의 송금자를 AWP와 송금 입차수보다 아래에 줄 세워요. 나중의 허락을 두 기간 모두에서 받은 허락 수보다 아래에, 그리고 W1에서는 활동 재시작을 쓴 엔도스랭크보다도 아래에 줄 세워요. 청산 없이 닫은 비율도 받은 허락 수보다 아래에 줄 세워요. 시빌 안정성 계수가 AWP보다 낮아요. 그리고 세 걷기 가운데 비용이 가장 많이 들어요.

결합 연산자(두 그래프의 걷기를 하나로 합치는 계산 규칙)는 $\lambda$의 정해 둔 값 세 개에서 시험했어요. 이때 고른 순간 이동을 쓰고, 점마다 이동 행렬 수준에서 볼록 혼합을 했어요. 두 층은 달러 금액으로 무게를 달았고, 각 허락은 $U$에서 잘랐어요. 엔도스랭크와 AWP처럼 무게를 단 층, 점에 따라 달라지거나 학습으로 정한 무게, 이동 행렬의 행 대신 정상 벡터를 섞는 방식, 그리고 다른 마르코프 연쇄(다음 걸음이 지금 있는 곳에만 달린 걷기)는 이 논문에서 평가하지 않아요. 그래서 RQ4와 RQ5의 답은 이 자료에서 그 방법 가족에만 들어맞아요. 실제로 쓰는 일은 여전히 \ref{sec:regulatory-ethics}절의 고려 사항에 묶여 있어요.

허락 쪽 비교 기준인 엔도스랭크는 허락 점검과 잘 맞고, AWP와 다르게 줄 세우고, 그래프가 듬성듬성해서(화살표가 적어서) 비용이 적게 들어요. 엔도스랭크는 송금 정보를 거의 담지 않아요. 그리고 두 기간 모두에서 나중의 허락을 원시 받은 허락 수보다 아래에 줄 세워요(라벨을 본 뒤에 계산했어요). 또 C-PR처럼 닫힌 농장에 드러나 있어요. 어떤 점수도 허락한 주인이 적은 계약이나, 점수를 매기지 않는 외부 소유 계정에 대해서는 아무것도 말해 주지 않아요(\ref{sec:data-sources}절).

자료는 네 가지 한계를 보태요. 그래프는 다섯 토큰만 담아요(표~\ref{tab:usd-tokens}). 그래서 주인들이 다른 토큰으로만 허락한 계약은 코호트 밖에 있어요. 모든 가격은 공개 서비스 하나, 곧 DefiLlama의 하루 가격 기록에서 와요 \parencite{defillamanddocs}. 그래서 어느 날짜의 가격에 잘못이 있으면, 그날의 모든 금액에 그 잘못이 옮겨 갈 거예요(부록~\ref{sec:data-tokens}). 계정 종류는 주소마다 한 번씩, 두 라벨 기간이 끝난 뒤인 2026년 10월 9일의 블록에서 읽었어요(표~\ref{tab:account-types}). 그래서 EIP-6780 전에 스스로를 없앤 계약은 외부 소유 계정으로 여겨져 빠져요(\ref{sub:cohorts}절) \parencite{ballet2023eip6780}. 그리고 트레이더 코호트는 작아요(\ref{sub:cohorts}절). 그래서 모든 트레이더 결과의 구간이 넓어요.

어떤 점수도 누가 주소를 움직이는지에 대한 오프체인(블록체인 밖의) 증거를 쓰지 않아요. 점수들의 오프체인 입력은 토큰 목록과 하루 가격 기록뿐이에요. 이 점수들은 휴먼 패스포트(Human Passport) 같은 증명 도구(사람인지 확인해 주는 도구)와 달라요 \parencite{humanpassportnddocs}. 이 점수들은 순위가 높은 두 주소를 서로 다른 주체가 움직인다고 보장하지 않아요. 그래서 이미 받아들인 무리 안의 순서로 읽어야 하고, 시빌 공격을 막는 장치로 읽으면 안 돼요. C-PR이나 엔도스랭크를 계산하기 전에, 증명 관문(통과해야 하는 확인 단계)을 통과한 주인의 허락만 세는 것은 그럴듯한 쓰임이에요. 하지만 관문 자체의 약점이 그 쓰임으로 옮겨 올 거예요 \parencite{bordeianu2026sybil}. 그리고 그 조합을 평가하려면 시빌인지 아닌지 표시한 라벨 자료가 필요한데, 이 연구는 그런 자료를 모으지 않았어요.

이 한계들은 내릴 수 있는 결론을 좁혀요. 이 계약 코호트와 이 기간에 대해, C-PR은 송금에 바탕을 둔 평판 점수에 허락 정보를 실어 넣는 걷기예요. 다만 들어오는 흐름에서는 손해를 봐요. 그리고 엔도스랭크는 스펜더 계약을 위한, 값싼 허락 전용 점수예요. 둘 다 신용 점수 방법이 아니고, 시빌 공격을 버티는 신원 장치도 아니에요. 서비스 수준의 담보를 대신하지도 못해요 \parencite{gudgeon2020crisis}. 이 점수들이 바탕을 둔 허락 화살표는 온체인 평판에 보탤 만한 가치가 있어요. 그리고 허락한 주인 수는 그것을 쓰는 가장 싼 방법이에요. \ref{ch:conclusion}장은 무엇을 보였는지, 그 앎이 어떤 가치가 있는지, 그리고 이 한계를 넓히려면 무엇이 필요한지 다시 정리해요.
